%File principale del documento su cui invocare la compilazione, vedi "istruzioni.txt" per più info

%Preambolo: la parte prima del \begin{document}
\documentclass[12pt,a4paper]{article} %formato del documento e grandezza caratteri

%Input del file metadata.tex della cartella locale "res/"
%lista di comandi presenti in template_latex.tex, da qui posso essere modificati secondo le esigenze

\newcommand{\DocTitle}{Verbale interno 2019-11-18} %variabile usata dal file template_latex.tex per settare il titolo del documento
%\newcommand{\DocAuthor}{Progetto "Predire in Grafana"} %variabile usata dal file template_latex.tex per settare l'autore del documento
\newcommand{\DocDate}{18 Novembre 2019} %variabile usata dal file template_latex.tex; Impostata manualmente, altrimenti ad ogni compilazione viene messa la data del giorno di compilazione.
\newcommand{\DocDesc}{Resoconto dell'incontro del gruppo \textit{VRAM Software} tenutosi in data 2019-11-18} %variabile usata dal file template_latex.tex per settare la descrizione del documento
\newcommand{\ver}{27.0.0} %variabile usata dal file template_latex.tex per settare la versione del documento
\newcommand{\app}{Toffoletto Massimo} %variabile usata dal file template_latex.tex per settare l'approvatore del documento
\newcommand{\red}{Dalla Libera Marco} %variabile usata dal file template_latex.tex per settare il redattore del documento
\newcommand{\test}{Schiavon Rebecca} %variabile usata dal file template_latex.tex per settare il verificatore del documento
\newcommand{\stat}{Approvato} %variabile usata dal file template_latex.tex per settare lo stato del documento
\newcommand{\use}{Interno} %variabile usata dal file template_latex.tex per indicare l'uso del documento %Contiene le varibili che descrivono il documento

%Input di file di configurazione presi dalla cartella "Template-LaTeX/config/", uguali per tutti i documenti
%Attenzione bisogna impostare il percorso del file!
% Tutti i pacchetti usati, da inserire nel preambolo prima delle configurazioni

\usepackage[T1]{fontenc} %Permette la sillabazione su qualsiasi testo contenente caratteri
\usepackage[utf8]{inputenc} %Serve per usare la codifica utf-8
\usepackage[english,italian]{babel} %Imposta italiano lingua principale, inglese secondaria. Es. serve per far apparire "indice" al posto di "contents"

\usepackage{graphicx} %Serve per includere le immagini

\usepackage[hypertexnames=false]{hyperref} %Gestisce i riferimenti/link. Es. Serve per rendere clickabili le sezioni dell'indice

\usepackage{float} %Serve per migliore la definizione di oggetti fluttuanti come figure e tabelle. Es. poter usare l'opzione [H] nelle figure ovvero tenere fissate le immagini che altrimenti LaTeX si sposta a piacere.

\usepackage{listings} %Serve per poter mettere snippets di codice nel testo

\usepackage{lastpage} %Serve per poter introdurre un'etichetta a cui si può fare riferimento Es. piè di pagina; poter fare " \rfoot{\thepage\ di \pageref{LastPage}} "

\usepackage{fancyhdr} %Per header e piè di pagina personalizzati

%Sono alcuni package che potranno esserci utili in futuro
%\usepackage{charter}
%\usepackage{eurosym}
\usepackage{subcaption}
%\usepackage{wrapfig}
%\usepackage{background}
\usepackage{longtable} % tabella che può continuare per più di una pagina
\usepackage[table]{xcolor} % ho dovuto aggiungere table in modo da poter colorare le row della tabella, dava: undefined control sequences
%\usepackage{colortbl}

\usepackage{dirtree} % usato per creare strutte tree-view in stile filesystem
\usepackage{xspace} % usato per inserire caratteri spazio
\usepackage[official]{eurosym}
\usepackage{pdflscape} %Inclusione pacchetti
% Configurazioni varie, da inserire nel preambolo dopo i pacchetti

\hypersetup{hidelinks} %serve per nascondere riquadri rossi che circondano i link 

\lstset{literate= {à}{{\`a}}1 } %Permette di usare lettere accentate nei listings

\pagestyle{fancy} %Imposto stile pagina
\fancyhf{} %Reset, se lo tolgo LaTex mette impostazioni di default (p.es numerazione pagine di default)


\lhead{\includegraphics[scale=0.25]{img/logo_header.png}} %Left header che compare in ogni pagina
%\rhead{\leftmark} %Nome della top-level structure (p.es. Section in article o Chapter in book) in ogni pagina
\rhead{\DocTitle\ v. \ver} %Right header

\newcommand{\glo}{$_G$} %Comando per aggiungere il pedice G
\newcommand{\glosp}{$_G$ } %Comando per aggiungere il pedice G con spazio

\newcommand\Tstrut{\rule{0pt}{2.6ex}} % top padding
\newcommand\Bstrut{\rule[-0.9ex]{0pt}{0pt}} % bottom padding
\newcommand{\TBstrut}{\Tstrut\Bstrut} % top & bottom padding

%Setto il colore dei link
%\hypersetup{
%	colorlinks,
%	linkcolor=[HTML]{404040},
%	citecolor={purple!50!black},
%	urlcolor={blue!50!black}
%}

%Tabelle e tabulazione (può tornare utile)
%\setlength{\tablcolsep}{10pt}
%\renewcommand{\arraystretch}{1.4}

%Comando per aggiungere le pagine di ogni sezione
%\newcommand{\newSection}[1]{%
%	\input{res/sections/#1}
%}

% Comandi per aggiungere padding a parole contenute nella tabella; è una specie di strut (un carattere invisibile)
%\newcommand\Tstrut{\rule{0pt}{2.6ex}} % top padding
%\newcommand\Bstrut{\rule[-0.9ex]{0pt}{0pt}} % bottom padding
%\newcommand{\TBstrut}{\Tstrut\Bstrut} % top & bottom padding  %Configurazione pacchetti

\begin{document}
	%Input del file "frontmatter" preso dalla cartella "Template-LaTeX/config/", uguale per tutti i documenti
	%Attenzione bisogna impostare il percorso del file!
	% #### FRONTESPIZIO (frontmatter) ####
\setlength{\headheight}{33pt} %Distanzia l'header
\pagenumbering{gobble} %Toglie il numero di pagina
\begin{titlepage}
	\begin{center}
		\includegraphics[scale=0.6]{img/logo.png} \\ %Logo
		\vspace{0.4cm} %Aggiunge uno spazio verticale di 0.5 cm
		
		{\LARGE Progetto "Predire in Grafana"} \\ %Nome progetto
		\vspace{0.4cm} %Attenzione a mettere il punto e NON la virgola
		
		{\Huge \textbf{\DocTitle}} \\ %Titolo, prende variabile definita in metadata.tex
		\vspace{0.4cm}
		
		\DocDate \\ %Data, prende variabile definita in metadata.tex
		\vspace{0.4cm}
		
		%Allineamento colonne: l=left r=right c=center, 
		%va specificato per ogni colonna
		%Se si vuole la riga tra colonne mettere "|"
		
		\begin{tabular}{r | l} %Elementi colonne separate da "&", le righe finiscono con "\\"
			Versione             & \ver \\
			Approvazione         & \app \\ 
			Redazione            & \red \\
			Verifica             & \test \\
			Stato                & \stat \\
			Uso                  & \use \\
		    Destinato a          & Zucchetti \\
						         & Prof. Vardanega Tullio\\
						         & Prof. Cardin Riccardo\\
			Email di riferimento & vram.software@gmail.com
		\end{tabular}
		\vfill
		\textbf{Descrizione} \\
		\DocDesc
	\end{center}
\end{titlepage}
\clearpage

% #### Impostazione header, footer  e numerazione pagine ####
\pagenumbering{arabic} %Pagine con i numeri arabi + reset a 1
\renewcommand{\footrulewidth}{0.4pt} %Di default footrulewidth==0 e quindi è invisibile, di default \headrulewith==0.4pt
\rfoot{\thepage\ di \pageref{LastPage}} %Pagina n di m, con numeri Arabi; usa il pacchetto "lastpage", in caso non sia possibile usare tale pacchetto mettere al fondo dell'ultima pagina "\label{LastPage}"

% #### Tabella dei log ####
% \textbf = grassetto; \Large = font più grande
% \rowcolors{quanti colori alternare}{colore numero riga pari}{colore numero riga dispari}: colori alternati per riga
% \rowcolor{color}: cambia colore di una riga
% p{larghezza colonna}: p è un tipo di colonna di testo verticalmente allineata sopra, ci sarebbe anche m che è centrata a metà ma non è precisa per questo utilizzo TBStrut; la sintassi >{\centering} indica che il contenuto della colonna dovrà essere centrato
% \TBstrut fa parte di alcuni comandi che ho inserito in config.tex che permetto di aggiungere un po' di padding al testo
% \\ [2mm] : questra scrittura indica che lo spazio dopo una break line deve essere di 2mm
% 

%\setcounter{secnumdepth}{0}
%\hfill \break
%\textbf{\Large{Diario delle modifiche}} \\


\addtocontents{toc}{\protect\setcounter{tocdepth}{0}} %Inserire questo per escludere una sezione dall'indice.

\section*{Registro delle modifiche} %Asterisco per fare sezione non numerata
\rowcolors{2}{gray!25}{gray!15}
\begin{longtable} {
		>{\centering}p{17mm} 
		>{\centering}p{19.5mm}
		>{\centering}p{24mm} 
		>{\centering}p{24mm} 
		>{}p{32mm}}
	\rowcolor{gray!50}
	\textbf{Versione} & \textbf{Data} & \textbf{Nominativo} & \textbf{Ruolo} & \textbf{Descrizione} \TBstrut \\
	14.7.0 & 2020-04-09 & Stantagiuliana Vittorio, Toffoletto Massimo e Spreafico Alessandro & \textit{Progettista}, \textit{Verificatore} e \textit{Responsabile di progetto} & Stesura, verifica e approvazione documento. \TBstrut \\ [2mm]
\end{longtable}

\addtocontents{toc}{\protect\setcounter{tocdepth}{4}} %Inserire questo per ripristinare il normale inserimento delle sezioni nell'indice. 4 significa fino al paragrah
\clearpage

% #### INDICE (tableofcontents) ####
\tableofcontents %Provoca la stampa dell'indice
\clearpage

\setcounter{secnumdepth}{4} %Permette di andare fino alla profondità del paragraph con la numerazione delle sezioni %Imposta il frontespizio, l'indice, header e footer
	
	\listoftables
	\clearpage
	
	\addtocontents{toc}{\protect\setcounter{tocdepth}{5}}
	\setcounter{secnumdepth}{5}
	
	
	%Tutte le sezioni del documento
	\section{Introduzione}
\subsection{Scopo del documento}
Il presente documento ha come scopo la definizione di regole e convenzioni che stanno alla base del progetto\glosp e che tutti i membri del gruppo devono seguire. In questo modo si garantisce consistenza e omogeneità in tutto il materiale del progetto\glosp stesso. Infatti ogni componente del gruppo è obbligato a visionare questo documento e a rispettarne rigorosamente le regole definite al fine di lavorare secondo una metodologia coesa e uniforme.
\subsection{Scopo del prodotto}
Il capitolato\glosp C4 ha lo scopo di implementare un plug-in di Grafana\glosp scritto in linguaggio Javascript che esegue gli algoritmi di Support Vector Machine SVM\glosp o Regressione Lineare RL\glo, i quali, leggendo da un file JSON la loro configurazione, saranno in grado di generare previsioni che potranno essere aggiunte al flusso di monitoraggio. Il plug-in prevede il superamento di determinati livelli di soglia per generare un allarme attraverso il meccanismo di Grafana\glosp detto alert\glo. Il plug-in utilizza un applicativo esterno per addestrare gli algoritmi di SVM\glosp ed RL\glosp o altri algoritmi di predizione.
Dunque il plug-in permette un'attività di monitoraggio grazie alla quale gli operatori possono intervenire con ogni cognizione di causa sul sistema.
\subsection{Ambiguità e glossario}
Questo documento è corredato da un \textit{Glossario v. 2.1.1} dove saranno illustrati i termini tecnici o altamente specifici per evitare ambiguità in essi. Le voci interessate sono identificate da una 'G' a pedice.
\subsection{Riferimenti}
\subsubsection{Riferimenti normativi}
\begin{enumerate}
	\item \textbf{Capitolato}\glosp \textbf{d'appalto C4 - Predire in Grafana}\glo: \url{https://www.math.unipd.it/~tullio/IS-1/2019/Progetto/C4.pdf};
\end{enumerate}
\subsubsection{Riferimenti informativi}
\begin{enumerate}
	\item \textbf{Terry Rout. Information technology - Software life cycle processes, Standard ISO/IEC 12207:1995 Australia\&New Zealand \\ Sezione: Primary life cycle processes. \\ Sezione: Supporting life cycle processes. \\ Sezione: Organizational life cycle processes}: \url{https://www.math.unipd.it/~tullio/IS-1/2009/Approfondimenti/ISO_12207-1995.pdf} (paragrafi §5.2, §5.3, §6.1-§6.5, §6.8, §7.1 e §7.4);
	\item \textbf{Slide L05 del corso Ingegneria del Software - Ciclo di vita del software}: \url{https://www.math.unipd.it/~tullio/IS-1/2016/Dispense/L05.pdf};
	\item \textbf{Slide L12 del corso Ingegneria del Software - Qualità di prodotto}: \url{https://www.math.unipd.it/~tullio/IS-1/2019/Dispense/L12.pdf}.
	\item \textbf{Standard ISO 25010}:
	\url{https://iso25000.com/index.php/en/iso-25000-standards/iso-25010}
	\item \textbf{Standard ISO 25023}:
	\url{https://metriche-per-il-software-pa.readthedocs.io/it/latest/documento-in-consultazione/metriche-e-strumenti.html#la-norma-iso-25023}
	
	%\item \textbf{Documentazione git}:
	%\url{https://www.atlassian.com/git};
	
	%\item \textbf{Standard ISO 8601}:
	%\url{https://it.wikipedia.org/wiki/ISO_8601}
	
	%\item \textbf{Snake Case}\glo:
	%\url{https://it.wikipedia.org/wiki/Snake_case} Pratica di scrivere tutto con gli underscore->va ancora deciso
	
	%\item \textbf{Airbnb JavaScript style guide}:
	%\url{https://github.com/airbnb/javascript/blob/master/README.md}
	
	%NOTA BENE: Dobbiamo decidere un po' di tecnologie da usare/strumenti e qualche sito per la loro documentazione da inserire nei riferimenti informativi
	
\end{enumerate}



	\pagebreak
	\section{Processi Primari}

\subsection{Fornitura}
\subsubsection{Scopo}
Lo scopo del processo di fornitura è di stabilire le norme e i tempi da rispettare nell'intero processo di sviluppo del progetto.
Andrà redatto uno Studio di Fattibilità al termine del quale sarà manifestata la decisione del capitolato scelto; sarà quindi necessario formalizzare un contratto per la consegna del prodotto insieme all'azienda proponente e stilare un Piano di Progetto. 
Per svolgere quanto sopra descritto sarà necessario seguire le seguenti fasi:
\begin{itemize}
	\item Avvio
	\item Contrattazione con il proponente
	\item Pianificazione del lavoro
	\item Esecuzione di quanto pianificato e continuo controllo della qualità
	\item Revisione e valutazione di quanto prodotto
	\item Consegna al proponente.
\end{itemize}
\subsubsection{Aspettative}
Il gruppo si aspetta di mantenere un dialogo costante con l'azienda proponente al fine di poter comprenderne al meglio le esigenze.
\subsubsection{Attività}
\paragraph{Studio di Fattibilità}
Lo Studio di Fattibilità viene redatto dagli analisti con lo scopo di stilare le motivazioni per cui il gruppo ha preferito un progetto rispetto agli altri.
Risulta quindi necessario individuare i costi e i benefici di ogni capitolato al fine di determinare quello che meglio si bilancia con le aspettative del gruppo tenendo sempre conto delle risorse a disposizione.
Per ogni capitolato andranno indicati:
\begin{itemize}
	\item \textbf{Informazioni Generali:} elenco che presenta le informazioni basilari, in particolare 
	\begin{itemize}
		\item Nome del capitolato
		\item Proponente
		\item Committente
	\end{itemize} 
	\item \textbf{Descrizione:} breve inquadramento dell'argomento del capitolato;
	\item \textbf{Prerequisiti di Progetto:} descrizione più dettagliata con un focus sugli obiettivi da soddisfare;
	\item \textbf{Vincoli:} elenco dei vincoli imposti dal proponente per la realizzazione del progetto;
	\item \textbf{Tecnologie Interessate:} elenco e sintetica descrizione delle tecnologie che dovranno essere impiegate durante lo svolgimento del capitolato;
	\item \textbf{Aspetti positivi:} elenco delle principali motivazioni che spingerebbero il gruppo a scegliere il capitolato;
	\item \textbf{Criticità e fattori di rischio:} elenco delle principali motivazioni che spingerebbero il gruppo a non scegliere il capitolato;
	\item \textbf{Conclusioni:} sintesi sulle motivazioni per cui il capitolato è stato scelto o è stato escluso.
\end{itemize}
\paragraph{Piano di Progetto}
Il Piano di Progetto viene redatto dal responsabile in collaborazione con gli amministratori e ha lo scopo di pianificare le attività per la realizzazione del progetto e individuare le risorse disponibili e da assegnare a ciascuna di esse.
Fanno parte di questo documento:
\begin{itemize}
	\item \textbf{Analisi dei rischi:} analisi dei rischi che potrebbero sorgere nel ciclo di sviluppo del progetto e delle modalità previste per prevenirli o affrontarli nel caso non sia stato possibile evitarli;
	\item \textbf{Modello di sviluppo:} descrizione del modello di sviluppo selezionato per lo svolgimento del progetto;
	\item \textbf{Pianificazione:} specifica dei tempi e dei ruoli per ogni attività individuata al fine di affrontare al meglio le scadenze del capitolato e la limitata quantità di risorse;
	\item \textbf{Preventivo e Consultivo:} stima dei costi necessari per lo svolgimento di ogni scadenza prevista e conseguente costruzione del preventivo per la totalità del progetto. Stesura di un verbale interno su cui riportare l'andamento effettivo rispetto a quello previsto.
\end{itemize}
\paragraph{Piano di Qualifica}
Il Piano di Qualifica viene redatto dai verificatori e ha lo scopo di descrivere le tecniche di verifica e validazione che i verificatori useranno nelle loro attività.
Questo documento presenterà:
\begin{itemize}
	\item \textbf{Qualità di Processo:} individuazione dei processi da adottare e delle metriche per misurarne la qualità;
	\item \textbf{Qualità di Prodotto:} individuazione delle caratteristiche del prodotto necessarie per una buona qualità e delle metriche per misurarne la qualità;
	\item \textbf{Specifiche dei test:} individuazione dei test a cui sottoporre il prodotto per garantire la soddisfazione dei suoi requisiti;
	\item \textbf{Standard di Qualità:} elenco degli standard di qualità scelti;
	\item \textbf{Valutazioni per il miglioramento:} esposizione dei problemi rilevati e delle soluzioni da attuare;
	\item \textbf{Conclusioni:} resoconti dei risultati delle misurazioni di qualità effettuate per ogni attività.
\end{itemize}

\subsubsection{Strumenti}
Durante il Processo di Fornitura sono utilizzati i seguenti strumenti della suite di Microsoft
\paragraph{Microsoft Excel}
Per svolgere semplici calcoli, creazione di grafici, diagrammi e tabelle si è scelto di utilizzare il software Microsoft Excel.
\paragraph{Microsoft Project}
Per la pianificazione, in particolare per la gestione di tempi e risorse e per l'analisi dei carichi di lavoro si è scelto di creare i diagrammi di Gantt e per farlo è stato utilizzato Microsoft Project.

//eventuali screen di Excel e/o Project

\subsection{Sviluppo}
\subsubsection{Scopo}
Lo scopo dello sviluppo è quello di svolgere le attività necessarie per la realizzazione del prodotto richiesto.
\subsubsection{Aspettative}
Il gruppo si aspetta di, fissati gli obiettivi e i vincoli, realizzare un prodotto completo e di buona qualità che sia ritenuto idoneo dal proponente
\subsubsection{Descrizione}
Lo sviluppo di divide nelle seguenti attività:
\begin{itemize}
	\item Analisi dei Requisiti
	\item Progettazione
	\item Codifica
\end{itemize}
\subsubsection{Attività}
\paragraph{Analisi dei Requisiti} \mbox{}
\\L'Analisi dei Requisiti viene redatta dagli analisti.
\\ \textbf{Scopo}\\ 
Lo scopo di questo documento è quello di individuare ed analizzare i requisiti del progetto al fine di:
\begin{itemize}
	\item individuare lo scopo del lavoro;
	\item delineare dei riferimenti precisi da fornire ai progettisti che saranno da seguire per l'intero processo di sviluppo e ai verificatori per le attività di controllo della qualità;
	\item creare dei riferimenti per la stima dei costi.
\end{itemize}
\textbf{Descrizione}\\
L'analisi dei requisiti viene redatta eseguendo le seguenti attività:
\begin{itemize}
	\item \textbf{comprensione del capitolato} attraverso la lettura e l'analisi della documentazione fornita;
	\item \textbf{comunicazioni con il proponente} al fine di migliorare l'analisi del progetto;
	\item \textbf{comunicazioni interne al gruppo}
	\item \textbf{analisi dei possibili casi d'uso}
\end{itemize}
\textbf{Casi d'uso}\\
Un caso d'uso definisce l'insieme di utenti, all'interno di uno scenario, con un obiettivo finale comune. Questo insieme è ottenuto attraverso la valutazione di ogni requisito.
Per descrivere ogni caso d'uso viene utilizzata la seguente struttura:
\begin{itemize}
	\item Codice Identificativo;
	\item Titolo;
	\item Diagramma UML;
	\item Attori primari;
	\item Attori secondari;
	\item Descrizione;
	\item Scenario;
	\item Scenario alternativo (ove presente);
	\item Inclusioni (ove presenti)
	\item Estensioni (ove presenti)
	\item Specializzazioni (ove presenti);
	\item Precondizione;
	\item Postcondizione;	
\end{itemize}
Il codice identificativo sarà scritto in questo formato: \\
\textbf{UC[codice\_padre].[codice\_figlio]} \\
Dove:
\begin{itemize}
	\item \textbf{codice\_padre:} numero che identifica univocamente i casi d'uso;
	\item \textbf{codice\_figlio:} numero progressivo che identifica i sottocasi;
\end{itemize}
\textbf{Requisiti}\\
Per descrivere un requisito viene utilizzata la seguente struttura:
\begin{itemize}
	\item Codice Identificativo;
	\item Classificazione;
	\item Descrizione;
	\item Fonti.
\end{itemize} 
Il \textbf{Codice Identificativo} sarà scritto in questo formato: \\
\textbf{R[Importanza][Tipologia][Codice]} \\
Dove:
\begin{itemize}
	\item \textbf{Importanza} può assumere i seguenti valori:
	\begin{itemize}
		\item 1: requisito obbligatorio;
		\item 2: requisito desiderabile;
		\item 3: requisito opzionale.
	\end{itemize}
	\item \textbf{Tipologia} può assumere i seguenti valori:
	\begin{itemize}
		\item F: Funzionale;
		\item Q: Prestazionale;
		\item P: Qualitativo;
		\item V: Vincolo.
	\end{itemize}
	\item\textbf{Codice}: numero progressivo identificativo strutturato nel formato: [codice\_padre].[codice\_figlio]
\end{itemize}
Le \textbf{Fonti} possono essere:
\begin{itemize}
	\item capitolato: il requisito è stato quindi individuato dalla lettura del capitolato;
	\item interno: il requisito è stato individuato ed aggiunto in seguito ad un'analisi interna;
	\item caso d'uso: il requisito è stato individuato dallo studio di un caso d'uso;
	\item proponente: il requisito è stato individuato in seguito ad un colloquio con il proponente.
\end{itemize}
\textbf{UML} \\
I diagrammi UML vengono realizzati usando la versione 2.0 del linguaggio.
\paragraph{Progettazione} \mbox{} \\
L'attività di progettazione è compito dei Progettisti.
\\ \textbf{Scopo}\\
Lo scopo della progettazione è di descrivere una soluzione soddisfacente per gli stakeholder. 
Viene quindi definita l'architettura del prodotto finale che dovrà:
\begin{itemize}
	\item Soddisfare i requisiti individuati nell'Analisi dei Requisiti;
	\item Essere comprensibile;
	\item Essere modulare;
	\item Garantire robustezza al fine di riuscire a gestire errori improvvisi.
\end{itemize}
\textbf{Descrizione}\\
La progettazione è composta da due sezioni:
\begin{itemize}
	\item Technology baseline;
	\item Product baseline.	
\end{itemize}
\textbf{Tecnology Baseline}\\
Nella Tecnology Baseline si trovano:
\begin{itemize}
	\item Descrizione delle \textbf{Tecnologie Utilizzate} con l'indicazione di vantaggi e svantaggi;
	\item Descrizione dei \textbf{Design pattern} utilizzati con corrispondente diagramma esplicativo;
	\item Elenco dei \textbf{Diagrammi UML}, in particolare:
	\begin{itemize}
		\item Diagrammi delle Classi;
		\item Diagrammi dei Package;
		\item Diagrammi di Attività;
		\item Diagrammi di Sequenza.
	\end{itemize}
	\item \textbf{Tracciamento delle Componenti:} riferimento per ogni requisito al componente atto a soddisfarlo;
	\item Descrizione dei \textbf{Test di Integrazione}, necessari per verificare che l'unione delle parti funzioni nel modo previsto.
\end{itemize}
\textbf{Product Baseline}\\
Nella Product Baseline si trovano:
\begin{itemize}
	\item \textbf{Descrizione delle Classi};
	\item \textbf{Tracciamento delle Classi:} riferimento per ogni requisito alla classe che lo soddisfa;
	\item \textbf{Test di Unità:} necessari per verificare che ogni parte funzioni individualmente nel modo previsto.
\end{itemize}
\paragraph{Codifica} \mbox{}
\\\textbf{Scopo}\\
Lo scopo di questa sezione è quello di normare tutte le regole di codifica al fine di garantire leggibilità e manutenibilità del codice. I Programmatori dovranno quindi seguire le seguenti regole.
\\\textbf{Stile di Codifica}\\
In seguito vengono elencate le norme da rispettare nel processo di codifica.
\begin{itemize}
	\item \textbf{Lingua:} la lingua con cui andranno scritti codice e commenti dovrà essere l'inglese;
	\item \textbf{Indentazione:} i blocchi annidati dovranno essere indentati e si dovranno utilizzare per ogni livello 4 spazi;
	\item \textbf{Nomi:} i nomi di classi, metodi e variabili dovranno essere univoci ed esplicativi;
	\item \textbf{Classi:} il nome di ciascuna classe dovrà iniziare con la lettera maiuscola; 
	\item \textbf{Costanti:} il nome di ciascuna costante dovrà essere scritto usando solo lettere maiuscole;
	\item \textbf{Metodi:} il nome di ciascun metodo deve iniziare con una lettera minuscola ad eccezione dei nomi composti che saranno scritti con le iniziali di ogni singola parola che li compone maiuscola 
	\item \textbf{Ricorsione:} l'uso della ricorsione è da evitare ove possibile.	
\end{itemize}
\subsubsection{Strumenti}
\paragraph{[nome strumento]}
	\pagebreak
	\section{Processi di supporto}
    \subsection{Documentazione}
        \subsubsection{Descrizione}
            In questo capitolo verrà illustrato come il gruppo ha gestito il processo di redazione della documentazione utile ai fini del progetto.
            Si andranno perciò a chiarire le fasi di costruzione del documento, la sua struttura, gli standard tipografici utilizzati e gli strumenti
            usati per la stesura.
        \subsubsection{Previsioni}
            Tale sezione si prefigge due obiettivi:
            \begin{itemize}
                \item standardizzare la scrittura della documentazione al fine di offrire degli elaborati consistenti e coerenti tra di loro;
                \item dare delle linee guida per la produzione di documenti corretti e validi.
            \end{itemize}
        \subsubsection{Obiettivi}
            Questo documento stabilisce delle norme per la produzione e la verifica dei prossimi elaborati. Tali indicazioni saranno condivise e seguite
            da tutti gli elementi del gruppo.
        \subsubsection{Ciclo di vita dei documenti}
            Ogni documento, durante il suo ciclo di vita, attraversa tre fasi principali:
            \begin{itemize}
                \item \textbf{Stesura}: in questo stadio viene effettuata la redazione del documento seguendo le norme indicate e i template disponibili.
                                        Il responsabile, a seguito di un'analisi preliminare, divide (se necessario) il lavoro tra vari redattori che procederanno
                                        alla compilazione.
                \item \textbf{Verifica}: terminata la fase di stesura i redattori sottometteranno l'elaborato al controllo dei verificatori. I verificatori, 
                                         nominati in precedenza dal responsabile, analizzeranno il documento. Nel caso in cui si evidenzino delle criticità
                                         segnaleranno le dovute modifiche ai redattori, in alternativa comunicheranno al responsabile la conformità del materiale.
                \item \textbf{Approvazione}: alla fine il documento viene vagliato dal responsabile che ne approva ufficialmente il rilascio.
            \end{itemize}
        \subsubsection{Struttura dei documenti}
            \paragraph{Template \LaTeX} \mbox{}\\ \mbox{}\\
                Per la produzione di documenti consistenti e coerenti tra loro il gruppo ha deciso di creare un template \LaTeX\glosp che permette, al momento della
                stesura, di concentrarsi unicamente sul contenuto e di non impegnare troppo tempo nella strutturazione della pagina. Il template è formato da
                più parti, tra cui:
                \begin{itemize}
                    \item un file per la corretta formattazione della prima pagina;
                    \item un file contenente i pacchetti \LaTeX\glosp utilizzati per una migliore rappresentazione dei contenuti;
                    \item un file per modificare i comandi personalizzati, i quali garantiscono una maggiore modularità del template;
                    \item due file per la gestione delle tabelle.
                \end{itemize}
            \paragraph{Prima pagina} \mbox{}\\ \mbox{}\\
                Il frontespizio è strutturato in questo modo:
                \begin{itemize}
                    \item \textbf{Logo del gruppo}: corredato dal nome, il tutto centrato orizzontalmente;
                    \item \textbf{Capitolato\glo}: nome del progetto scelto dal gruppo (\textit{Predire in Grafana}), centrato orizzontalmente;
                    \item \textbf{Titolo del documento}: nome esplicativo del contenuto del documento, centrato orizzontalmente e in grassetto;
                    \item \textbf{Data}: data di approvazione finale dell'elaborato, centrata orizzontalmente;
                    \item \textbf{Tabella informativa}: contenente informazioni quali:
                    \begin{itemize}
                        \item versione del documento;
                        \item nominativo dell'approvatore (e quindi responsabile) del documento;
                        \item nominativo/i del/i redattore/i del documento;
                        \item nominativo/i del/i verificatore/i del documento;
                        \item stato del documento (approvato o non approvato);
                        \item tipo d'uso del documento (interno o esterno);
                        \item nominativi dei destinatari del documento;
                        \item email di riferimento del gruppo.
                    \end{itemize} 
                    \item \textbf{Descrizione}: breve riassunto del contenuto del documento.
                \end{itemize}
            \paragraph{Registro delle modifiche} \mbox{}\\ \mbox{}\\
                La pagina successiva mostra una tabella contenente una raccolta cronologica delle modifiche avvenute sul documento.
                Le colonne espongono i seguenti contenuti:
                \begin{itemize}
                    \item versione del documento dopo la modifica;
                    \item data della modifica;
                    \item nominativo del componente responsabile della modifica;
                    \item ruolo del componente responsabile della modifica;
                    \item breve descrizione della modifica effettuata;
                \end{itemize}
            \paragraph{Indice} \mbox{}\\ \mbox{}\\
                L'indice riassume in modo schematico i contenuti del documento mostrando la loro suddivisione in sezioni, sottosezioni e paragrafi.
                È presente all'inizio di ogni elaborato dopo il registro delle modifiche. I suoi elementi sono cliccabili e rimandano alla pagina
                corrispondente. In alcuni documenti viene corredato di due liste:
                \begin{itemize}
                    \item \textbf{Lista delle immagini}: riportante le immagini utilizzate tramite rispettiva didascalia;
                    \item \textbf{Lista delle tabelle}: riportante le tabelle utilizzate tramite rispettiva didascalia (sono esclusi il registro delle modifiche
                                                  e il riepilogo tracciamenti)-
                \end{itemize}
            \paragraph{Contenuto} \mbox{}\\ \mbox{}\\
                Il corpo principale del documento è così strutturato:
                \begin{itemize}
                    \item intestazione in alto con logo del gruppo riadattato a sinistra e nome del documento a destra;
                    \item una riga divisoria;
                    \item il contenuto della pagina formattato con una dimensione di 12 punti e la visualizzazione in A4 formato articolo
                          (\verb|\documentclass{article}|);
                    \item una riga divisoria opzionale in caso di appunti;
                    \item eventuali note sul testo riportate con un numero progressivo a esponente;
                    \item una riga divisoria;
                    \item il contenuto a piè di pagina con a destra il numero progressivo della pagina visualizzata sul totale (nell'indice invece dei numeri
                          arabi vengono utilizzati quelli romani).
                \end{itemize}
        \subsubsection{Norme tipografiche}
                \paragraph{Nomi dei file} \mbox{}\\ \mbox{}\\
                    Per organizzare i file in modo più ordinato e corretto il gruppo ha scelto di seguire la pratica dello "snake\_case"\glo.
                    Tale prassi, nell'ambito di questo progetto, è applicata seguendo tre regole:
                    \begin{itemize}
                        \item i file con nome composto da più parole useranno il carattere underscore come separatore;
                        \item i nomi dei file saranno scritti interamente in minuscolo;
                        \item non vengono tralasciate preposizioni di alcun tipo dal nome del file.
                    \end{itemize}
                \paragraph{Glossario}
                    \begin{itemize}
                        \item Ogni parola da inserire nel \textit{Glossario} è corredata da una "G" a pedice per tutte le sue occorrenze;
                        \item sono escluse dalla regola precedente i titoli e il nome delle sezioni del documento.
                    \end{itemize}
                \paragraph{Stili di testo}
                \begin{itemize}
                    \item \textbf{Grassetto}: il grassetto viene usato solamente per evidenziare le parole a cui si vuole dare risalto e per la prima parola
                                              di un elenco puntato se quest'ultimo è del tipo: \textbf{Termine}: definizione/spiegazione.
                    \item \textbf{Corsivo}: il corsivo viene usato per evidenziare i nomi dei documenti oltre che per il nome dell'azienda
                                            proponente (\textit{Zucchetti}) e del gruppo (\textit{VRAM Software}).
                    \item \textbf{Maiuscolo}: in aggiunta alle normali regole dell'italiano le lettere maiuscole verranno utilizzate anche per riportare i nomi dei
                                              documenti e di certi capitoli.
                \end{itemize}
                \paragraph{Elenchi puntati} \mbox{}\\ \mbox{}\\
                    Ogni voce dell'elenco puntato sarà preceduta da un "•" (punto elenco) per quanto riguarda il primo livello e da un "-" per il secondo.
                    Alla fine di ogni elemento, invece, si aggiunge il carattere ";" a meno che non sia l'ultimo dove va messo il segno di punteggiatura ".".
                    Se la voce dell'elenco è del tipo termine-spiegazione/definizione il termine va in grassetto.
                \paragraph{Convenzioni} \mbox{}\\ \mbox{}\\
                    Vengono usate tre convenzioni per le date, la valuta e gli orari (i quali prendono come riferimento ISO 8601). Le date saranno espresse nella forma: \newline \newline
                    \centerline{\textbf{DD-MM-YYYY}} \newline \newline
                    dove:
                    \begin{itemize}
                        \item \textbf{DD}: indica il giorno e deve sempre avere due cifre.
                        \item \textbf{MM}: indica il mese e deve sempre avere due cifre;
                        \item \textbf{YYYY}: indica l'anno;
                    \end{itemize}
                    Gli orari saranno invece rappresentati nella forma: \newline \newline
                    \centerline{\textbf{HH:MM}} \newline \newline
                    dove:
                    \begin{itemize}
                        \item \textbf{HH}: indica le ore (in formato 24 ore) e deve sempre avere due cifre;
                        \item \textbf{MM}: indica i minuti e deve sempre avere due cifre.
                    \end{itemize}
                    La valuta infine sarà scritta nella forma: \newline \newline
                    \centerline{\textbf{X.XXX,YY}} \newline \newline
                    dove:
                    \begin{itemize}
                        \item \textbf{XXXX}: rappresenta la parte intera della somma;
                        \item \textbf{.}: viene usato come separatore ogni tre cifre intere;
                        \item \textbf{YY}: rappresenta la parte decimale della somma;
                    \end{itemize}
            \subsubsection{Produzione}
                    Il progetto prevede la stesura di documenti che si dividono in interni ed esterni a seconda dei destinatari. Questi elaborati verranno poi valutati
                    in quattro revisioni. Tali elementi progettuali saranno elencati e spiegati di seguito con le rispettive sigle.
                \paragraph{Documenti esterni}
                    \begin{itemize}
                        \item \textbf{Analisi dei Requisiti - AdR}: contiene la descrizione degli attori del sistema e la loro interazione con i casi d'uso,
                                                                    fornendo una visione chiara ai progettisti sul problema da trattare;
                        \item \textbf{Manuale Utente - MU}: serve per guidare l'utente nell'utilizzo del software;
                        \item \textbf{Manuale Sviluppatore - MS}: ha lo scopo di spiegare a eventuali sviluppatori esterni il funzionamento
                                                                  del sistema;
                        \item \textbf{Piano di Progetto - PdP}: dà indicazioni sulle tempistiche e i costi che il gruppo prevede d'impiegare
                                                                per la realizzazione del progetto;
                        \item \textbf{Piano di Qualifica - PdQ}: fornisce le specifiche riguardanti il controllo qualità dei processi e dei
                                                                 prodotti tramite delle metriche misurabili.
                    \end{itemize}
                \paragraph{Documenti interni}
                    \begin{itemize}
                        \item \textbf{Glossario - G}: documento dove sono riportate i termini che necessitano di spiegazione o che potrebbero causare confusione o
                                                      fraintendimenti;
                        \item \textbf{Norme di Progetto - NdP}: linee guida per la gestione delle attività di progetto;
                        \item \textbf{Studio di Fattibilità - SdF}: breve analisi di tutti i capitolati proposti.
                    \end{itemize}
                \paragraph{Verbale} \mbox{}\\ \mbox{}\\
                    Il verbale della riunione è un documento particolare in quanto può essere sia interno che esterno. Sarà interno quando
                    i partecipanti saranno solo i componenti del gruppo; esterno quando saranno presenti anche referenti dell'azienda.
                    È inoltre presente la tabella "Riepilogo dei tracciamenti" che tiene segno delle modifiche discusse e
                    approvate in sede d'incontro. Questo prospetto presenta una colonna con un codice nel formato: \newline \newline
                    \centerline{\textbf{V(I/E)\_X.Y}} \newline \newline
                    dove:
                    \begin{itemize}
                        \item \textbf{V(I/E)}: indica se il verbale è interno o esterno;
                        \item \textbf{X}: indica il numero dell'incontro che si sta organizzando;
                        \item \textbf{Y}: indica il numero della scelta.
                    \end{itemize}
                \paragraph{Revisioni}
                \begin{itemize}
                    \item \textbf{Revisione dei Requisiti - RR}: analisi iniziale del capitolato per concordare i requisiti col proponente;
                    \item \textbf{Revisione di Progettazione - RP}: revisione intermedia per accertare la fattibilità del progetto;
                    \item \textbf{Revisione di Qualifica - RQ}: revisione intermedia per l'approvazione delle verifiche e l'attivazione del processo di validazione;
                    \item \textbf{Revisione di Accettazione - RA}: revisione finale per il collaudo del software e l'accertamento del soddisfacimento dei requisiti
                                                                   presentati in RR.
                \end{itemize}
                \paragraph{Altre sigle}
                    \begin{itemize}
                        \item \textbf{Responsabile - Re};
                        \item \textbf{Amministratore - Am};
                        \item \textbf{Analista - An};
                        \item \textbf{Progettista - Pr};
                        \item \textbf{Programmatore - Pt};
                        \item \textbf{Verificatore - Ve};
                        \item \textbf{Proof of Concept - PoC}.
                    \end{itemize}
            \subsubsection{Elementi grafici}
                    \paragraph{Immagini} \mbox{}\\ \mbox{}\\
                        Le immagini sono sempre centrate e presentano una didascalia, tranne il logo nella prima pagina.
                    \paragraph{Tabelle} \mbox{}\\ \mbox{}\\
                        Le tabelle sono sempre centrate e presentano una didascalia, tranne il "Registro delle modifiche" e il "Riepilogo tracciamenti".
                        Inoltre le tabelle presentano un'alternanza di colori per favorirne la lettura.
                    \paragraph{Didascalie} \mbox{}\\ \mbox{}\\
                        Le didascalie vengono indicate con un numero progressivo del tipo X.Y, dove:
                        \begin{itemize}
                            \item \textbf{X}: indica la sezione a cui appartiene l'elemento;
                            \item \textbf{Y}: indica che tabella o immagine è della sezione.
                        \end{itemize}
            \subsubsection{Strumenti}
                    \paragraph{\LaTeX} \mbox{}\\ \mbox{}\\
                        Per stilare i documenti è stato usato il linguaggio di markup \LaTeX\glo. Questo linguaggio permette di scrivere dei testi mantenendo il focus
                        sul contenuto e non sulla forma. La struttura infatti è stata dichiarata all'inizio del progetto e poi riutilizzata permettendo una gestione
                        degli elaborati scalabile, modulare e ordinata.
                    \paragraph{TeXStudio} \mbox{}\\ \mbox{}\\
                        Per scrivere in \LaTeX\glosp viene usata l'IDE \textit{TeXStudio}\glosp creata appositamente per questo fine e che quindi offre delle scorciatoie
                        per i vari comandi, un correttore automatico per la lingua italiana e un servizio di compilazione e visualizzazione dei PDF integrato;
                    \paragraph{Lucidchart} \mbox{}\\ \mbox{}\\
                        Per la creazione dei diagrammi UML\glosp e di altro tipo viene usata la piattaforma online Lucidchart\glo. \newline \newline
                        \centerline{\url{https://www.lucidchart.com}}
	\pagebreak
	\subsection{Gestione della configurazione}
	\subsubsection{Scopo}
		L'obiettivo di questo processo è rendere documenti e codice sorgente univocamente identificati e facilmente riconoscibili, evidenziando versioni e modifiche. Vuole anche agevolare l'identificazione delle relazioni esistenti fra gli elementi e fa da supporto alla fase di verifica.
	\subsubsection{Aspettative}
		Questa sezione ha lo scopo supportare i processi\glosp di documentazione, di sviluppo e manutenzione del software rendendoli definiti e ripetibili.  
	\subsubsection{Descrizione} 
		Questo processo descrive come il gruppo ha gestito la configurazione degli strumenti e delle risorse utilizzate per svolgere il progetto\glo.
		Sarà quindi descritta la configurazione del repository\glosp su GitHub, del sistema di versionamento\glosp Git e dei servizi GitHub.
	\subsubsection{Attività}
		\paragraph{Versionamento}\mbox{}\\ [1mm]
		Le modifiche effettuate ai documenti ed ai file contenenti codice sorgente sono identificate da un numero di versione presente all'interno dei file stessi. Per quanto riguarda i documenti, ogni numero di versione sarà presente nel nome del file ed avrà una corrispondente riga nella tabella delle modifiche così da avere uno storico delle modifiche effettuate al documento. Per il codice sorgente del software le modifiche effettuate in una versione saranno consultabili sulla sezione rilasci di GitHub e saranno identificate da opportuni tag.
		\newline
		La numerazione per il versionamento del codice sorgente si applicherà in modo indipendente alle singole componenti software che compongono il prodotto finale, tuttavia esisterà un'unico versione per il prodotto finale, che seguirà le linee guida del versionamento del codice sorgente basate su risoluzione di problemi e nuove funzionalità minori/maggiori. Nelle note di rilascio del prodotto finale saranno poi riportate le versioni dei singoli componenti che lo compongono ovvero documentazione, plugin Grafana\glosp e applicativo di addestramento esterno. Le modifche ed i rilasci saranno tracciabili tramite il repository principale grafana\_prediction che contiene al suo interno gli altri repository come sottomoduli.
		\subparagraph*{Numerazione delle versioni del documento}
		Per identificare univocamente le versioni dei documenti seguiamo le indicazioni dell'ente ETSI, dunque la numerazione delle versioni è composta da 3 cifre nel formato X.Y.Z e la prima bozza ha versione 0.0.1. Le modifiche al documento aggiorneranno le cifre nel seguente modo:
		\begin{itemize}
			\item \textbf{Z} : Questa cifra viene incrementata ogni volta che vengono effettuate delle modifiche editoriali\glosp al documento, es. X.Y.1, X.Y.2 etc.;
			\item \textbf{Y} : Questa cifra viene incrementata ogni volta che vengono effettuate delle modifiche tecniche\glosp al documento. Se sono state eseguite modifiche sia editoriali\glosp che tecniche\glosp allora entrambe le cifre Z e Y saranno incrementate, es. X.1.1, X.2.2 etc;
			\item \textbf{X} : Questa cifra viene incrementata ogni volta che viene rilasciata una nuova versione finale del documento, ovvero una versione di rilascio.			
		\end{itemize}
		La prima versione finale ha versione 1.1.1.
		\newline
		Durante la modifica e la revisione del documento le bozze vengono incrementate le cifre Z ed Y, es. 1.2.0, 1.2.1, 1.3.0, etc.
		\newline
		Quando il documento viene approvato come finale allora viene incrementata la cifra X, ad esempio la bozza 1.5.3 diventa la versione finale 2.1.1.
		\newline
		Ogni incremento sulle singole cifre è rigorosamente un +1, non possono essere saltati numeri.
		
		\subparagraph*{Numerazione delle versioni del codice sorgente}
		Per identificare univocamente le versioni del codice sorgente utilizziamo una numerazione basata su 3 cifre nel formato X.Y.Z, la prima bozza avrà versione 0.0.0 e le modifiche aggiorneranno le cifre della versione nel modo seguente:
		\begin{itemize}
			\item \textbf{Z} : Questa cifra viene incrementata ogni volta che vengono effettuate delle modifiche al codice sorgente al fine di risolvere problemi riscontrati nel software;
			\item \textbf{Y} : Questa cifra viene incrementata ogni volta che vengono effettuate delle modifiche al codice sorgente al fine di aggiungere nuove funzionalità di entità minore al software, la nuova versione deve inoltre presentare piena compatibilità con le versioni precedenti. Se sono state eseguite modifiche sia per risolvere problemi che per aggiungere nuove funzionalità minori allora entrambe le cifre Z e Y saranno incrementate, es. X.1.1, X.2.2 etc;
			\item \textbf{X} : Questa cifra viene incrementata ogni volta che vengono effettuate modifiche al codice al fine di aggiungere nuove funzionalità di entità maggiore al software, ad esempio novità che cambino in modo importante l'uso del software o le funzionalità che esso offre. Un aumento di versione di questo tipo può, se necessario, non essere retrocompatibile.
		\end{itemize}
		La prima versione finale ha versione 1.0.0.
		\newline
		Ogni incremento sulle singole cifre è rigorosamente un +1, non possono essere saltati numeri.
		
	\paragraph{Gestione delle modifiche}\mbox{}\\ [1mm]
		Al fine di monitorare e limitare le modifiche al ramo principale del repository\glo, master, è utilizzato il meccanismo di pull request fornito da GitHub. Ogni membro del gruppo può creare branch secondari, secondo il workflow\glosp feature branch, su cui effettuare modifiche, tuttavia per unirle al branch master è necessario aprire una pull request che dovrà essere revisionata dai verificatori tramite i servizi di revisione integrati in GitHub. Una volta revisionata positivamente è compito del responsabile del documento approvare la pull request ed effettuare quindi l'effettiva unione delle modifiche nel branch master.
		\newline
		In sintesi, per effettuare modifiche ai file sono previsti i seguenti passaggi:
		\begin{itemize}
			\item contattare il responsabile del file affinché autorizzi la modifica del file stesso;
			\item creare un branch secondario ed effettuare le modifiche al file;
			\item aprire una pull request per unire le modifiche al ramo master;
			\item i verificatori revisionano la pull request ed eventualmente richiedono aggiornamenti;
			\item completata la revisione il responsabile approva la pull request.
		\end{itemize}
	
		Il meccanismo appena esposto di applica anche al codice sorgente, dove inoltre sono presenti dei sistemi di verifica automatica e di continuous integration che possono bloccare le pull request se le modifiche effettuate non rispettano i livelli di qualità desiderati o se interrompono la compilazione o il superamento dei test automatici sul codice. In particolare, gli strumenti utilizzati a questo scopo sono le Github Actions per implementare la continuous integration, JEst per l'esecuzione dei test automatici sul codice sorgente ed i servizi Coveralls e SonarCloud per controllare la qualità del codice sorgente.
		
	\paragraph{Repository}\mbox{}\\ [1mm]
		Per tenere traccia di versioni e modifiche fatte a documenti e codice è utilizzato il sistema di versionamento\glosp distribuito Git, che può essere utilizzato tramite riga di comando o utilità grafiche come GitHub Desktop o GitKraken.
		La struttura dei repository utlizzati è la seguente:
		\newline
		\dirtree{%
			.1 grafana\_prediction.
			.2 VRAM\_SW\_DOC\_2020.
			.2 grafana\_prediction\_plugin.
			.2 prediction\_configuration.
		}
		\mbox{}\\ % forza un newline
		Dove grafana\_prediction è il repository principale e gli altri sono suoi sottomoduli, configurati tramite la funzionalità Git submodules di Git.
		Il repository\glosp principale, contenente quindi i sottomoduli, è ospitato sul sito GitHub all'indirizzo: 
		\begin{center}
			\url{https://github.com/VRAM-Software/grafana_prediction}
		\end{center}
		\subparagraph{Struttura del repository VRAM\_SW\_DOC\_2020}
		Al fine di fornire una navigazione agevole e standardizzata, il contenuto del repository\glosp è organizzato in modo gerarchico tramite directory secondo il seguente schema:
		\newline
		\dirtree{%
			.1 root.
				.2 RR.
					.3 Esterni.
					.3 Interni.
				.2 RP.
					.3 Esterni.
					.3 Interni.
				.2 RQ.
					.3 Esterni.
					.3 Interni.
				.2 RA.
					.3 Esterni.
					.3 Interni.
				.2 Guide.
				.2 Template.
					.3 Images.
					.3 config.
					.3 img.
					.3 res.
		}
		\mbox{}\\ % forza un newline
		Nel dettaglio:
		\begin{itemize}
			\item \textbf{RR}: contiene i sorgenti \LaTeX \xspace dei documenti relativi alla revisione dei requisiti;
			\item \textbf{RP}: contiene i sorgenti \LaTeX \xspace dei documenti relativi alla revisione di progettazione\glo;
			\item \textbf{RQ}: contiene i sorgenti \LaTeX \xspace dei documenti relativi alla revisione di qualifica;
			\item \textbf{RA}: contiene i sorgenti \LaTeX \xspace dei documenti relativi alla revisione di accettazione; 
			\item \textbf{Guide}: contiene brevi indicazioni interne su come usare i comandi \LaTeX \xspace usati nei documenti e su come riutilizzare il template per generare nuovi documenti;
			\item \textbf{Template}: contiene i sorgenti \LaTeX \xspace usati per generare la base comune di tutti i documenti.
		\end{itemize}
		È inoltre disponibile una directory condivisa su Google Drive, con la stessa struttura gerarchica sopra esposta, che contiene tutti e soli i file PDF in versione finale. Lo scopo di questa directory condivisa è facilitare la consultazione e la condivisione dei file PDF stessi.
		
	\paragraph{Tipologie di file non accettate}\mbox{}\\ [1mm]
		Tramite un apposito file .gitignore presente nella root directory della gerarchia vengono definiti i tipi di file non accettati all'interno del repository\glo. Vengono esclusi tutti i file di compilazione, compilati o temporanei in quanto il repository\glosp dovrebbe contenere solamente i seguenti formati di file:
		\begin{itemize}
			\item file sorgente \LaTeX \xspace con estensione .tex;
			\item file immagine, preferibilmente in formato .png;
			\item file testuali in formato .md o .txt;
			\item tutti i file di codice sorgente come .js, .ts etc.		
		\end{itemize}
	\subsubsection{Strumenti di supporto}
	\paragraph{GitHub}
	\subparagraph*{Strumenti GitHub utilizzati}\mbox{}\\ [1mm]
		In aggiunta ai servizi già elencati, al fine di migliorare efficacia ed efficienza, vengono utilizzate le funzionalità di "Issue Tracking System",
		% milestone qui è usato come nome di funzione github quindi non va a glossario  
		"Milestone" e "Project Board" integrate in GitHub. Ognuna di queste funzionalità viene usata solo da chi autorizzato, ad esempio rilasci di versioni, creazione e chiusura di milestone sono concesse solo al responsabile di progetto\glo.
		\newline
		Vengono inoltre utilizzate le "Labels" offerte dallo Issue Tracking System Github di per definire le tipologie di problemi e le loro priorità.
		\newline
		Per il repository\glosp contenente il codice sorgente del software sono utilizzate le GitHub Actions al fine di implementare la pratica della continuous integration, che utilizzerà strumenti quali JEst, SonarCloud e Coveralls per eseguire le verifiche automatiche sul codice sorgente.
	\paragraph{SonarCloud}
	I repository del codice sorgente sono stati collegati al servizio SonarCloud ed è stata inserita la sua analisi statica del codice sorgente nei controlli eseguiti dalla continuous integration. Le pagine sonarcloud dei repository sono disponibili al seguente indirizzo: 
	\begin{center}
		\url{https://sonarcloud.io/organizations/vram-software/projects}
	\end{center}
	\paragraph{Coveralls}
	Poiché SonarCloud, nella versione disponibile al momento dello svolgimento del progetto\glo, non supporta l'analisi del code coverage tramite Github Actions utilizziamo il servizio offerto da Coveralls e sempre integrato nella continuous integration. Le pagine Coveralls dei repository sono disponibili ai seguenti indirizzi:
	\begin{center}
		\url{https://coveralls.io/github/VRAM-Software/grafana_prediction_plugin}
	\end{center}
	\begin{center}
	\url{https://coveralls.io/github/VRAM-Software/prediction_configuration_utility}
	\end{center}
	
		
	\pagebreak
	\subsection{Gestione della qualità}
\subsubsection{Scopo}
Lo scopo della gestione della qualità è quello di dare una valutazione oggettiva riguardo ai servizi che il prodotto\glosp deve fornire in modo mantenere uno standard di qualità prestabilito e da soddisfare i requisiti impliciti ed espliciti del proponente.
\subsubsection{Aspettative}
Le aspettative della gestione della qualità possono essere riassunti dai seguenti punti, che rappresentano gli obiettivi che vogliono essere raggiunti dal gruppo. Si vuole ottenere:
\begin{itemize}
	\item qualità per tutta la durata del ciclo di vita del software;
	\item una copertura totale per tutti i requisiti impliciti ed espliciti richiesti dal proponente;
	\item la soddisfazione del cliente con il prodotto\glosp sviluppato;
	\item oggettività nella ricerca della qualità in modo da ottenerla in tutti i processi e in tutti i prodotti\glo;
\end{itemize}
\subsubsection{Descrizione}
Per la gestione della qualità dello sviluppo di questo progetto si è deciso di cercare di mettere in atto l'idea dello standard ISO 9000, il quale afferma che il sistema di qualità, cioè l'insieme dei metodi e regole per gestire la qualità di un prodotto\glo, è costituito da tre passaggi: lo sviluppo, il controllo e il miglioramento continuo. Questi tre passaggi sono ripetuti per ogni processo in modo di cercare di prevenire possibili errori e anomalie.
Per una descrizione approfondita sul piano della qualità, che è uno dei costituitenti del sistema di qualità, si faccia riferimento al seguente documento: \textit{Piano di Qualità v. 1.1.1}.
Nel documento sopracitato sono descritte le metodologie e regole utilizzate dal gruppo \textit{VRAM Software} per cercare di raggiungere le aspettative definite nella sezione precedente.
Il contenuto sintetico del documento \textit{Piano di Qualità v. 1.1.1} contiene:
\begin{itemize}
	\item la presentazione delle caratteristiche e gli attributi considerati più importanti del prodotto\glo;
	\item la descrizione degli standard di qualità utilizzati per la gestione del processo di vita del software;
	\item l'individuazione dei processi definiti dagli standard descritti, nei quali sono successivamente stati individuati gli obiettivi, le strategie per raggiungerli e le metriche per misurarne l'avanzamento.
\end{itemize}
Con il controllo e il miglioramento dei processi a ogni prodotto\glosp vengono poi associati: gli obiettivi da raggiungere e le metriche utilizzate.
L'obiettivo che si vuole raggiungere, utilizzando le regole proposte nel documento \textit{Piano di Qualità v. 1.1.1}, è quello di ottenere un livello di qualità soddisfacente per software e documentazione, perchè se si riesce a perseguire la qualità per tutta il ciclo di vita del prodotto\glo, allora le possibilità di ottenere un prodotto\glosp stabile, fonte di miglioramento continuo, sarà più alta.
\subsubsection{Attività}
Per ogni processo abbiamo deciso di pianificare lo sviluppo in tre attività:
\begin{itemize}
	\item \textbf{Pianificazione}
		\begin{itemize}
			\item ci siamo posti degli obiettivi per perseguire la qualità;
			\item abbiamo definito delle strategie per raggiungere tali obiettivi;
			\item abbiamo disposto le persone e le risorse da utilizzare per la realizzazione del processo.
		\end{itemize}
	\item \textbf{Valutazione} \\*
		Abbiamo messo in atto le strategie proposte nella pianificazione valutandone i risultati;
	\item \textbf{Reazione} \\*
		Per ogni risultato ottenuto abbiamo modificato opportunamente le nostre strategie.
\end{itemize}
\subsubsection{Strumenti}
Gli strumenti utilizzati per la gestione della qualità sono:
\begin{itemize}
	\item strumenti forniti dallo standard ISO-12207, che contiene la descrizione del processo: "Gestione della Qualità";
	\item le metriche utilizzate per misurare il perseguimento della qualità.
\end{itemize}

	\pagebreak
	\subsection{PRC-Q6 Verifica}
\subsubsection{Scopo}
Il processo\glosp di verifica ha lo scopo di accertare che i risultati ottenuti da tutte le attività, attuate in un periodo preso in considerazione, raggiungano tutti i requisiti richiesti senza aver introdotto degli errori. La verifica deve essere svolta durante il ciclo di vita del software, più precisamente al passare di ogni baseline\glo.
\subsubsection{Aspettative}
La verifica viene effettuata per raggiungere i seguenti obiettivi:
\begin{itemize}
	\item cercare consistenza, completezza e correttezza nelle singole attività;
	\item ottenere supporto per una successiva validazione\glosp del prodotto\glo;
	\item avere a disposizione dei criteri e procedure, comprensibili e affidabili, da seguire.
\end{itemize}
\subsubsection{Descrizione}
Il processo\glosp di verifica assicura che tutte le attività svolte della fase in esame, attraverso analisi e test successivamente definiti, siano conformi alle aspettative.
\subsubsection{Attività}
Le attività che si effettuano durante la verifica sono due tipi di analisi: l'analisi statica\glosp e l'analisi dinamica. \\[2mm]
\paragraph{Analisi statica}\mbox{}\\ [1mm]
L'analisi statica\glosp è lo studio della documentazione e del codice sorgente, essa non richiede l'esecuzione delle singole parti del sistema software che si sta sviluppando, quindi viene applicata ad ogni prodotto\glosp di processo\glo.
\subparagraph*{Verifica della documentazione} \mbox{} \\ [1mm] 
La verifica della documentazione può essere effettuata manualmente in due modi:
\begin{itemize}
	\item \textbf{Walkthrough}: metodo di lettura in cui tutti i componenti del gruppo analizzano tutti i documenti e codice realizzato, senza tralasciare nulla, in modo da rilevare eventuali errori. Questa ricerca è ad ampio spettro, infatti viene effettuata anche se a priori non si può sapere se verranno trovate anomalie;
	\item \textbf{Inspection}: metodo di lettura o desk check per cercare errori specifici utilizzando liste di controllo o check list, la tabella successiva rappresenta gli errori che i nostri verificatori hanno identificato come essere i più frequenti.
\end{itemize}
\rowcolors{2}{gray!25}{gray!15}
\setcounter{table}{0}
\begin{longtable} {
		>{}p{55mm} 
		>{}p{55mm}
		>{}p{0mm}}
	\rowcolor{gray!50}
	\textbf{Oggetto da controllare} & \textbf{Correzione} & \TBstrut \\[2mm]
	Formato data dei documenti & Tutte le date devono utilizzare la seguente scrittura: \textit{YYYY-MM-DD}  &  \TBstrut \\[2mm]
	Sintassi nome documenti citati & I documenti, interni o esterni, citati nella documentazione devono essere scritti in italico utilizzando il comando: \verb|\textit{}|  &  \TBstrut \\[2mm]
	Sintassi parole da enfatizzare & Le parole da enfatizzare, come titolo di liste puntate, devono essere scritte i grassetto utilizzando il comando: \verb|\textbf{}|  &  \TBstrut \\[2mm]
	Punteggiatura liste puntate & Ogni elemento nelle liste puntate deve iniziare con una lettera maiuscola o minuscola (a seconda del contesto) e deve finire con un punto e virgola &  \TBstrut \\[2mm]
	Uso corretto dei comandi \verb|\glo| e \verb|\glosp| & I comandi \verb|\glo| e \verb|\glosp|, sono stati creati per identificare le parole da inserire nel glossario, \verb|\glo| dovrà essere usato se la parola del glossario precede un carattere di punteggiatura. Il comando \verb|\glosp| sarà invece usato per le parole che precedono uno spazio &  \TBstrut \\	[2mm]
	\rowcolor{white}
	\caption{Tabella errori più frequenti}
\end{longtable}

\paragraph{Analisi dinamica} \mbox{}\\ [1mm]
L'analisi dinamica è una tecnica che consiste nell'esecuzione del prodotto\glosp software eseguendo dei test su un insieme finito di casi.
\subparagraph{Test}\mbox{}\\ [1mm]
Per la corretta verifica del codice, i test devono seguire regole ben precise che definiscono i parametri e le proprietà che devono avere test per essere considerati efficaci: \mbox{}\\ [1mm]
\textbf{Parametri da definire}:
\begin{itemize}
	\item \textbf{Input}: l'input preso in considerazione;
	\item \textbf{Output}: l'output che si aspetta di ricevere;
	\item \textbf{Ambiente}: l'ambiente, quindi hardware e sistema operativo, in cui viene eseguito il test;
	\item \textbf{Stato iniziale}: lo stato iniziale del prodotto\glo, precedente all'esecuzione del test;
	\item \textbf{Istruzioni opzionali}: l'utilizzo di istruzioni opzionali aggiuntive deve essere noto.
\end{itemize} 
Un test può essere definito efficace se:
\begin{itemize}
	\item i parametri citati sopra sono correttamente indicati;
	\item il nome di un test deve essere autoesplicativo per capire subito che cosa si sta testando;
	\item esso è ripetibile, quindi ci si aspetta un risultato costante per tutte le esecuzioni di quel test;
	\item in caso di fallimento del test, esso devo fornire informazioni utili a risolvere il problema identificato;
	\item in caso di possibili effetti indesiderati ci si aspetta un avvertimento;
	\item esso è esaustivo e accurato, infatti devono poter identificare una parte del progetto\glosp che potrebbe essere causa di errori.
\end{itemize}
Il software testing è un procedimento importante per misurare la qualità di un prodotto\glosp software, per questo motivo sono definiti quattro livelli che dividono i tipi di test da effettuare secondo una gerarchia. Nel processo\glosp di verifica essi sono: 
\begin{itemize}
	\item \textbf{Test di sistema};
	\item \textbf{Test di integrazione};
	\item \textbf{Test di unità}.
\end{itemize}
\subparagraph*{Codifica dei test} \mbox{} \\ [1mm]
Per ogni test vengono definiti il codice identificativo, la descrizione, lo stato e l'esito.
\begin{itemize}
	\item il codice identificativo è scritto in questo formato:\\
	\\ \textbf{tipologia[codice\_padre].[codice\_figlio]} \\
	\\ Dove:
	\begin{itemize}
		\item \textbf{tipologia}:
		\begin{itemize}
			\item \textbf{TS}: test di sistema;
			\item \textbf{TI}: test di integrazione;
			\item \textbf{TU}: test di unità.
		\end{itemize}
		\item \textbf{codice\_padre}: numero intero progressivo maggiore di zero che identifica univocamente il test;
		\item \textbf{codice\_figlio}: numero intero progressivo maggiore di zero che indica i sotto-test che compongono il padre. 
	\end{itemize}
	\item descrizione: spiegazione concisa e completa di ciò che verifica il test;
	\item stato:
	\begin{itemize}
		\item \textbf{I}: implementato;
		\item \textbf{NI}: non implementato.
	\end{itemize}
	\item Esito:
	\begin{itemize}
		\item \textbf{P}: positivo;
		\item \textbf{N}: negativo;
		\item \textbf{NE}: non eseguito.
	\end{itemize}
\end{itemize}

\paragraph*{Test di unità} \mbox{}\\ [1mm]
I test di unità si definiscono tali perché controllano il corretto funzionamento di una sezione di codice specifica chiamata unità. Solitamente questi tipi di test vengono eseguiti su funzioni o, se si parla di programmazione ad oggetti, su un metodo di una classe, quindi su una parte minimale del codice. Il loro scopo è verificare che gli input forniti diano gli output attesi. È buona norma rendere questi test automatici per minimizzare il tempo di esecuzione senza richiedere l'interazione dello sviluppatore. Questo tipo di test viene scritto dal programmatore che sviluppa le singole unità mediante l'ausilio di driver\glosp e stub\glo, per verificare l'assenza di errori e documentare il funzionamento dell'unità.

\paragraph*{Test di integrazione} \mbox{}\\ [1mm]
Successivamente ai test di unità, si assemblano gruppi di unità progressivamente che formeranno degli agglomerati sempre più grandi. Questo tipo di test, infatti, si concentra nella ricerca di anomalie ed errori tra le interfacce dei singoli componenti. Una volta concluso un singolo test, si otterrà un agglomerato che costituirà una nuova unità da cui partire per ulteriori test di integrazione. Per svolgere questo test si è scelto un approccio iterativo.

\paragraph*{Test di sistema} \mbox{}\\ [1mm]
I test di sistema controllano il funzionamento corretto del prodotto\glosp e verificano che il sistema soddisfi tutti i requisiti definiti nel documento \textit{Analisi dei Requisiti v. 3.1.1}. Con l'esecuzione di questi test si riesce a definire e documentare tutte le funzionalità del sistema e possono individuare criticità come il problema del black-box\glo. Il gruppo svolgerà questo test quando si raggiungerà uno stato di rilascio per il prodotto\glosp sviluppato.

\subsubsection{Obiettivi di qualità}
	\paragraph{OP-4 Efficacia dei test}\mbox{} \\ [1mm]
	%\paragraph*{Obiettivo}\mbox{}\\ [1mm]
	L'obiettivo è garantire che il prodotto\glosp software venga testato in modo efficace in ogni sua singola parte.
	\subparagraph{Metriche}\mbox{}\\ [1mm]
	Le metriche\glosp utilizzate sono:
	\begin{itemize}
		\item \textbf{M-PROC05 Percentuale bug sistemati}: indica la percentuale dei bug sistemati dando così una traccia dell'avanzamento;
		\begin{itemize}
			\item[] \textbf{formula}: $\frac{bug \; sistemati}{bug \; totali}$;
		\end{itemize}
	\end{itemize}

	\paragraph{OP-5 Completezza dei test} \mbox{} \\ [1mm]
	%\paragraph*{Obiettivo}\mbox{}\\ [1mm]
	L'obiettivo è garantire che il prodotto software sia testato in modo completo.
	\subparagraph{Metriche}\mbox{}\\ [1mm]
	Le metriche\glosp utilizzate sono:
	\begin{itemize}
		\item \textbf{M? Code coverage}: indica la percentuale di codice coperto da test rispetto al totale delle righe di codice;
		\begin{itemize}
			\item[] \textbf{formula}: $\frac{righe \; di \; codice \; coperto}{righe \; di \; codice \; totali}$;
		\end{itemize} 
	
		\item \textbf{M? Branch coverage}: indica la percentuale di codice coperto da test rispetto al totale delle righe di codice;
		\begin{itemize}
			\item[] \textbf{formula}: $\frac{righe \; di \; codice \; coperto}{righe \; di \; codice \; totali}$;
		\end{itemize} 
	
		\item \textbf{M? Copertura dei test eseguiti}: indica la percentuale di test effettivamente eseguiti;
		\begin{itemize}
			\item[] \textbf{formula}: $\frac{test \; eseguiti}{test \; totali}$;
		\end{itemize} 
	\end{itemize}
	\paragraph{PRD-Q2 Appropriatezza funzionale}
		\subparagraph{CP-3 Correttezza delle funzionalità implementate} \mbox{}\\ [1mm]
		%\paragraph*{Caratteristica}\mbox{}\\ [1mm]
		Questa caratteristica descrive se le funzionalità implementate sono corrette.
		\subparagraph{Metriche}\mbox{}\\ [1mm]
		Le metriche\glosp utilizzate sono:
		\begin{itemize}
			\item \textbf{M? percentuale di test superati}: indica la percentuale di test che sono stati superati;
			\begin{itemize}
				\item[] \textbf{formula}: $\frac{test \; superati}{test \; eseguiti}$;
			\end{itemize}
		\end{itemize}
	
\subsubsection{Strumenti di supporto}
	\paragraph{Correzione ortografica}\mbox{}\\ [1mm]
	Il nostro gruppo ha utilizzato, per la ricerca di errori ortografici o stilistici nel testo dei documenti, la libreria HunSpell utilizzata con l'editor di test: TexStudio. Esso permette di avere un controllo in tempo reale su errori ortografici in lingua inglese e italiana (sottolineati in rosso) e ripetizioni di parole a breve distanza (sottolineati in verde).
	\paragraph{SonarJS}\mbox{}\\ [1mm]
	Per l'analisi statica\glosp del codice Javascript viene utilizzato il code analyzer SonarJS\glosp della piattaforma SonarCloud.
	\paragraph{SonarTS}\mbox{}\\ [1mm]
	Per l'analisi statica\glosp del codice Typescript viene utilizzato il code analyzer SonarJS\glosp della piattaforma SonarCloud.
	\paragraph{Validazione W3C}\mbox{}\\ [1mm]
	Le pagine di markup scritte il linguaggio HTML5 vengono validate attraverso il validatore W3C reperible al seguente indirizzo: 
	\begin{center}
		\url{https://validator.w3.org/}
	\end{center}
	Le pagine di presentazione scritte in linguaggio CSS3 vengono validate attraverso il validatore W3C reperibile al seguente indirizzo:
	\begin{center}
		\url{https://jigsaw.w3.org/css-validator/}
	\end{center}
	\paragraph{Jest}\mbox{}\\ [1mm]
	Per la gestione dei test del codice JavaScript il nostro gruppo ha utilizzato il framework Jest per eseguire test d'unità sul codice React e NodeJS.
	\paragraph{Enzyme}\mbox{}\\ [1mm]
	Abbiamo scelto di utilizzare Enzyme in quanto è uno strumento di supporto specifico per la scrittura di test sui componenti React e si integra con il framework Jest.
	\paragraph{React Testing Library}\mbox{}\\ [1mm]
	A supporto del framework Jest, viene utilizzata la libreria Testing Library per React che facilita la scrittura di test manutenibili per componenti React.
	\paragraph{Coveralls}\mbox{}\\ [1mm]
	Coveralls è uno strumento che permette di quantificare la percentuale di codice testato.

	\pagebreak
	\subsection{Validazione}

\subsubsection{Scopo}
Lo scopo del processo\glosp di validazione\glosp è accertare che il prodotto\glo, per uno specifico uso, corrisponda alle attese.

\subsubsection{Aspettative}
Le aspettative del gruppo sono di normare adeguatamente il processo\glosp di validazione\glosp al fine di rilasciare un prodotto\glosp privo di errori gravi e capace di rispondere alle richieste del proponente.

\subsubsection{Descrizione}
Le attese sono quelle del committente e corrispondono, se stilata correttamente, all'\textit{Analisi dei Requisiti v. 3.1.1}. 
Per eseguire una validazione\glosp si deve avere un prodotto\glosp che è a un sufficiente livello di avanzamento, quindi il numero di validazioni svolte è molto minore rispetto al numero delle verifiche, che iniziano a venire eseguite prima.
Durante questo processo vengono svolti i test di accettazione tramite i quali il proponente valida il prodotto\glosp software.

\subsubsection{Attività}
	\paragraph{Test di accettazione}
		\subparagraph*{Codifica del test}
		Per ogni test vengono definiti il codice identificativo, la descrizione, lo stato e l'esito.
		\begin{itemize}
			\item il codice identificativo è scritto in questo formato:\\
			\textbf{tipologia[codice\_padre].[codice\_figlio]}
			\begin{itemize}
				\item \textbf{tipologia}:
				\begin{itemize}
					\item \textbf{TA}: test di accettazione;
				\end{itemize}
				\item \textbf{codice\_padre}: numero che identifica univocamente il test;
				\item \textbf{codice\_figlio}: numero progressivo che indica i sotto-test che compongono il padre. 
			\end{itemize}
			\item descrizione: spiegazione concisa e completa di ciò che verifica il test;
			\item stato:
			\begin{itemize}
				\item \textbf{I}: implementato;
				\item \textbf{NI}: non implementato.
			\end{itemize}
			\item Esito:
			\begin{itemize}
				\item \textbf{P}: positivo;
				\item \textbf{N}: negativo;
				\item \textbf{NE}: non eseguito.
			\end{itemize}
		\end{itemize}
		\subparagraph*{Descrizione} \mbox{}\\ [1mm]
		I test di accettazione, o di collaudo, vengono effettuati ad ogni pre-rilascio, quindi dopo l'esecuzione dei test di sistema. Essi servono a confermare il soddisfacimento dei requisiti da parte del committente; infatti questi tipi di test ne richiedono la presenza. Se questi test vengono superati significa che il prodotto\glosp software è pronto per il rilascio. 
		Per identificare un test di accettazione viene utilizzata una codifica con il seguente formato:\\
		\textbf{TA[codice]}
		Dove codice è un numero progressivo nel formato: [X].[Y] con X e Y numeri maggiori di zero.
\subsection{Strumenti di supporto}
\paragraph{Jest}\mbox{}\\ [1mm]
Per la gestione dei test del codice JavaScript il nostro gruppo ha utilizzato il framework Jest per eseguire test d'unità sul codice React e NodeJS.
\paragraph{Enzyme}\mbox{}\\ [1mm]
In particolare, per quanto riguarda i test in React, abbiamo scelto Enzyme: uno strumento specifico per effettuare test sui componenti React che va ad integrarsi con Jest.
\paragraph{React Testing Library}\mbox{}\\ [1mm]
A supporto del framework Jest, viene utilizzata la libreria Testing Library per React che facilita la scrittura di test manutenibili per componenti React.
\paragraph{Coveralls}\mbox{}\\ [1mm]
Coveralls è uno strumento che permette di quantificare la percentuale di codice testato.
	\pagebreak
	\subsection{PRC-Q8 Gestione dei cambiamenti}
\subsubsection{Scopo}
Lo scopo di questo processo\glosp è fornire un metodo tempestivo, disciplinato e documentato per assicurare che tutti i cambiamenti, compresi i problemi, riscontrati nel corso del progetto\glo, siano identificati, analizzati e risolti.
\subsubsection{Aspettative}
L'applicazione di questo processo\glosp deve comportare maggior tempestività ed efficacia nell'individuazione e nella risoluzione dei problemi, con un conseguente aumento di efficacia ed efficienza del processo\glosp di verifica. Inoltre questo processo\glosp deve permettere di individuare i problemi maggiormente ricorrenti.
\subsubsection{Descrizione} 
Il processo\glosp di gestione dei cambiamenti analizza e risolve tutti i cambiamenti programmati e improvvisi, quali i problemi, riscontrati durante l'esecuzione dei processi\glosp di sviluppo e di fornitura, indipendentemente dalla loro natura e dalla loro origine. 
\subsubsection{Attività}
\paragraph{Implementazione del processo}\mbox{}\\ [1mm]
Deve essere istanziato il processo\glosp di gestione dei cambiamenti per gestire ogni nuova modifica, comprese le non conformità ed i problemi non pianificati, rilevato nei prodotti\glosp software o nelle attività. Questo processo\glosp deve essere conforme alle seguenti proprietà:
\begin{itemize}
	\item il processo\glosp deve essere circolare e chiuso, assicurando che:
	\begin{itemize} 
		\item tutti i cambiamenti siano prontamente segnalati e gestiti tramite il processo di risoluzione dei cambiamenti;
		\item i cambiamenti vengano presi in carico e gestiti;
		\item vengano inviate delle notifiche per informare gli interessati della presenza dell'eventuale problema;
		\item le cause del problema vengano identificate, analizzate e, dove possibile, eliminate;
		\item la risoluzione e le decisioni prese siano archiviate e storicizzate;
		\item lo stato del cambiamento sia tracciato, aggiornato e comporti delle notifiche al cambio di stato;
		\item venga mantenuto un registro di tutti i problemi riscontrati.
	\end{itemize}
	\item Il processo\glosp definisce uno schema per categorizzare e dare priorità ai cambiamenti. Questo schema è composto da:
	\begin{itemize} 
		\item \textbf{Identificativo}: codice numerico progressivo per identificare il singolo cambiamento;
		\item \textbf{Tipologia}: denota il tipo di cambiamento e può essere:
		\begin{itemize}
			\item \textbf{Editoriale}: cambiamento che comporta modifiche editoriali\glosp sulla documentazione;
			\item \textbf{Tecnico}: cambiamento che comporta modifiche tecniche\glosp sulla documentazione;
			\item \textbf{Funzionale}: cambiamento che influenza le caratteristiche funzionali del prodotto\glosp software;
			\item \textbf{Conformità}: cambiamento che comporta violazioni di conformità del prodotto\glosp software;
			\item \textbf{Validazione}: cambiamento riscontrato durante il processo\glosp di Validazione\glosp del prodotto\glosp software.
		\end{itemize}
		\item \textbf{Priorità}: può essere:
		\begin{itemize}
			\item bloccante;
			\item urgente;
			\item alta;
			\item media;
			\item bassa.
		\end{itemize}
		\item \textbf{Stato}: può essere:
		\begin{itemize}
			\item da fare;
			\item in corso;
			\item completato.
		\end{itemize}
	\end{itemize}
	\item I cambiamenti rilevati ed archiviati vengono analizzati per identificare quelli più frequenti ed eventuali tendenze;
	\item La modalità di gestione e di risoluzione individuate per i cambiamenti e per i problemi sono analizzate e valutate al fine di verificare che siano stati effettivamente risolti e che le modifiche siano state correttamente implementate nei prodotti\glosp software e nelle attività. Inoltre l'analisi deve permettere di individuare dei pattern di risoluzione ed aiutare a determinare quando vengono introdotti ulteriori problemi.			
\end{itemize}

\paragraph{Risoluzione dei problemi}\mbox{}\\ [1mm]
Quando dei cambiamenti, comprese le non conformità ed i problemi, sono rilevati in un prodotto\glosp software o in un'attività, deve essere realizzato un report che li descriva. Questo report deve essere usato come componente del ciclo chiuso del punto §3.6.4.1: dall'individuazione dei cambiamenti al processo di indagine, analisi e risoluzione, fino all'individuazione della causa ed alla sua analisi per individuare le tendenze.
Per svolgere quanto descritto, ci si appoggia all'Issue Tracking System di GitHub. Esso consente di aggiungere delle issues che corrispondono ai cambiamenti o ai problemi che si devono svolgere e di assegnarle ad incaricati e responsabili. Dunque viene assegnata un'istanza dello schema precedentemente descritto e, in base alle priorità, si procederà alla loro risoluzione.

\subsubsection{Metriche di qualità}
	\paragraph{OP-6 Risoluzione efficace dei problemi}
	\paragraph*{Obiettivo}\mbox{}\\ [1mm]
	L'obiettivo è garantire una corretta risoluzione dei problemi nel più breve tempo possibile.
	\paragraph*{Metriche}\mbox{}\\ [1mm]
	Le metriche\glosp utilizzate sono:
	\begin{itemize}
		\item \textbf{M-PROC06 Tempo medio di risoluzione degli errori}: indica il tempo medio di risoluzione degli errori segnalati. Se elevato, questo valore indica che la gestione dei cambiamenti non è adeguata;
		\begin{itemize}
			\item[] \textbf{formula}: $\frac{TRE}{NE}$ dove $TRE$ indica il tempo totale per la risoluzione degli errori e $NE$ il numero di errori risolti.
		\end{itemize} 
	\end{itemize}		

	\pagebreak
	\section{Processi organizzativi}
	\subsection{PRC-Q9 Gestione organizzativa}
		\subsubsection{Scopo}
			Lo scopo del processo\glosp di gestione organizzativa è definito dai seguenti punti:
			\begin{itemize}
				\item creazione di un modello organizzativo che specifica i rischi che possono verificarsi;
				\item definizione un modello di sviluppo da seguire;
				\item pianificazione del lavoro in base alle scadenze fissate;
				\item creazione e calcolo di un piano economico suddiviso tra i ruoli;
				\item definizione finale del bilancio sul totale delle spese.
			\end{itemize}
			Queste attività sono definite a cura del responsabile di progetto\glosp e redatte all'interno del documento \textit{Piano di Progetto v. 6.1.1}.
		\subsubsection{Aspettative}
			Gli obiettivi del processo\glosp di gestione organizzativa sono:
			\begin{itemize}
				\item coordinare i componenti del gruppo attraverso l'assegnazione di ruoli e compiti;
				\item coordinare la comunicazione tra i membri del gruppo con l'ausilio di strumenti che la rendano facile ed efficace;
				\item pianificare l'esecuzione delle attività in modo ragionevole;
				\item gestire le attività anche dal punto di vista economico;
				\item monitorare costantemente il team, i processi\glosp e i prodotti\glosp in modo efficace.
			\end{itemize}
		\subsubsection{Descrizione}
			Il processo\glosp di gestione organizzativa è composto dalle seguenti attività:
			\begin{itemize}
				\item definizione dello scopo;
				\item generazione delle istanze dei processi\glo;
				\item stima e pianificazione di risorse, costi e tempo per lo svolgimento del progetto\glo;
				\item assegnazione dei ruoli e, per ognuno di essi, delle attività;
				\item gestione del controllo e dell'esecuzione di tutte le attività;
				\item valutazione periodica dello stato e dell'esecuzione delle attività rispetto a quanto pianificato.
			\end{itemize}			
		\subsubsection{Attività}
		\paragraph{Pianificazione delle risorse}\mbox{}\\ [1mm]
		Per gestire le risorse disponibili per il nostro progetto\glosp dobbiamo monitorare il rispetto dei costi e dei tempi preventivati. A tal proposito definiamo i ruoli di progetto e il modo in cui vengono assegnati.
		\paragraph{Ruoli di progetto}
		
			\paragraph*{Assegnazione}\mbox{}\\ [1mm]
				I ruoli scelti corrispondono alla rispettiva figura aziendale e valgono per la durata di una milestone. Al termine di ognuna di esse, viene eseguita una rotazione dei ruoli così da permettere a ciascuno dei membri del gruppo di ricoprirli tutti almeno una volta.
				I ruoli previsti sono descritti di seguito.
			\paragraph*{Responsabile di progetto}\mbox{}\\ [1mm]
				La figura del responsabile di progetto\glosp è molto importante in quanto convergono su di esso le responsabilità di gestione, pianificazione, coordinamento e controllo dell'intero progetto\glo.
				Le sue mansioni sono:
				\begin{itemize}
					\item redazione dell'organigramma e del documento \textit{Piano di Progetto};
					\item controllo e coordinamento di risorse, componenti e attività del gruppo;
					\item relazione con il controllo di qualità interno al progetto\glo;
					\item analisi e gestione degli elementi più critici e i rischi;
					\item intermediazione nelle comunicazioni esterne ovvero con proponente e committente;
					\item elaborazione ed emissione di piani e scadenza;
					\item approvazione finale dei documenti.
				\end{itemize}
			\paragraph*{Amministratore di progetto}\mbox{}\\ [1mm]
				La figura dell'amministratore è responsabile dell'efficienza e dell'operatività dell'ambiente di sviluppo e della redazione e attuazione di piani e procedure di gestione per la qualità.
				Le sue mansioni sono:
				\begin{itemize}
					\item controllo versioni e configurazioni del prodotto\glo;
					\item gestione di tutta la documentazione del progetto\glo;
					\item direzione delle infrastrutture di supporto;
					\item risoluzione dei problemi sorti dalla gestione dei processi\glo;
					\item stesura del documento \textit{Norme di Progetto} per conto del responsabile di progetto\glo;
					\item collaborazione con il responsabile di progetto\glo alla reazione del documento \textit{Piano di Progetto}.
				\end{itemize}
			\paragraph*{Analista}\mbox{}\\ [1mm]
				La figura dell'analista è responsabile di tutta l'attività di analisi dei problemi e del dominio applicativo.
				Le sue mansioni sono:
				\begin{itemize}
					\item studio del dominio del problema;
					\item analisi della complessità, della fattibilità e dei requisiti del problema;
					\item stesura del documento \textit{Studio di Fattibilità};
					\item stesura del documento \textit{Analisi dei Requisiti}.
				\end{itemize}
			\paragraph*{Progettista}\mbox{}\\ [1mm]
				La figura del progettista è responsabile delle attività di progettazione\glo. Deve perciò gestire gli aspetti tecnologici e tecnici del progetto\glo.
				Le sue mansioni sono:
				\begin{itemize}
					\item redigere le specifiche tecniche del prodotto\glosp effettuando scelte il più efficienti e performanti possibili, sulla base delle tecnologie presenti;
					\item sviluppare l'architettura del prodotto\glosp tale da fornire una base stabile a manutenibile che soddisfi le specifiche tecniche.
				\end{itemize}
			\paragraph*{Programmatore}\mbox{}\\ [1mm]
				La figura del programmatore è responsabile dell'attività di codifica mirata alla realizzazione del prodotto\glosp compreso delle componenti necessarie allo svolgimento di test di verifica e validazione\glo.
				Le sue mansioni sono:
				\begin{itemize}
					\item eseguire l'attività di codifica sulla base del lavoro svolto da progettista;
					\item fornire dei componenti che permettano si svolgere test e validazione\glosp sul prodotto\glo.
				\end{itemize}
			\paragraph*{Verificatore}\mbox{}\\ [1mm]
				La figura del verificatore è responsabile di tutta l'attività di verifica.
				Si affida al documento \textit{Norme di Progetto} è definito anche lo standard da seguire per questa attività.
				Le sue mansioni sono:
				\begin{itemize}
					\item redazione della parte retrospettiva del \textit{Piano di Qualifica};
					\item ispezione dei documenti al fine di controllare che siano conformi allo standard previsto nelle \textit{Norme di Progetto};
					\item segnalazione di eventuali errori nelle parti di prodotto\glosp a chi ha responsabilità su di esse.
				\end{itemize}
			\paragraph{Esecuzione e controllo}
				Il compito di eseguire e monitorare la gestione organizzativa viene assegnato all'amministratore. Egli infatti deve analizzare e risolvere i problemi che sorgono, fornendo dei report sullo stato di avanzamento di progetto sia internamente che al committente.	
		\paragraph{Gestione dei rischi}\mbox{}\\ [1mm]
		La gestione e l'analisi dei rischi è un'attività a carico del responsabile di progetto\glosp che deve essere documentata nel documento \textit{Piano di Progetto v. 6.1.1}.
		La procedura per la gestione dei rischi è composta dai seguenti passi:
		\begin{itemize}
			\item individuazione dei rischi;
			\item analisi dei rischi;
			\item pianificazione dell'attività di controllo;
			\item controllo e monitoraggio.
		\end{itemize}
		Viene definita una codifica per l'identificazione dei fattori di rischio: R[tipo][X]
		\begin{itemize}
			\item tipo:
			\begin{itemize}
				\item \textbf{RT}: rischi tecnologici di lavoro e produzione del software;
				\item \textbf{RG}: rischi di gruppo;
				\item \textbf{RO}: rischi organizzativi delle attività da svolgere;
				\item \textbf{RR}: rischi legati ai requisiti e ai rapporti con gli stakeholder\glo;
				\item \textbf{RS}: rischi di stima di costi e tempi.
			\end{itemize}
			\item X: numero intero progressivo che inizia da 1.
		\end{itemize}
		Inoltre vengono definiti delle misure di probabilità e degli indici di gravità:
		\begin{itemize}
			\item probabilità:
			\begin{itemize}
				\item \textbf{Alta};
				\item \textbf{Media};
				\item \textbf{Bassa}.
			\end{itemize}
			\item Gravità:
			\begin{itemize}
				\item \textbf{1}: basso livello di gravità;
				\item \textbf{2}: medio livello di gravità;
				\item \textbf{3}: alto livello di gravità.
			\end{itemize}
		\end{itemize}
		\subsubsection{Procedure}
		Sono state definite un insieme di procedure che definiscono la gestione del coordinamento interno del gruppo.
		\paragraph{Gestione delle comunicazioni}
			\subparagraph{Comunicazioni interne}\mbox{}\\ [1mm]
				Le comunicazioni interne al gruppo avvengono tramite la piattaforma Slack nella quale è stato creato un workspace suddiviso nei seguenti canali tematici:
				\begin{itemize}
					\item \textbf{Documentazione}: canale dedicato alle discussioni in merito alle attività di stesura e verifica della documentazione;
					\item \textbf{Sviluppo}: canale dedicato alle discussioni in merito alle attività di analisi, progettazione\glosp e codifica del prodotto\glo;
					\item \textbf{Generale}: canale dedicato alle comunicazioni di servizio per il coordinamento di attività, incontri e impegni del gruppo.
				\end{itemize}
				Nei canali dedicati alla documentazione e allo sviluppo è stata attivata l'integrazione con GitHub per segnalare i nuovi commit e le pull request con un messaggio.
				Inoltre è stata eseguita l'integrazione con Google Calendar che permette di avere una sezione dedicata alle notifiche degli eventi programmati sul calendario e l'integrazione del plug-in Simple Poll che permette di eseguire dei sondaggi per prendere le decisioni.
			\subparagraph{Comunicazioni esterne}\mbox{}\\ [1mm]
				Le comunicazioni esterne avvengono tramite la posta elettronica all'indirizzo \url{vram.software@gmail.com}.
				Il responsabile di progetto\glosp ha l'incarico di gestire queste comunicazioni ed è tenuto ad informare tutti gli elementi del gruppo degli argomenti trattati qualora non fossero presenti.
		\paragraph{Gestione degli incontri}
			\subparagraph{Incontri interni}\mbox{}\\ [1mm]
				Gli incontri del gruppo sono accordati su Slack nel canale generale oppure nella precedente riunione e vengono annotati su Google Calendar. Essi sono identificati da:
				\begin{itemize}
					\item data;
					\item ora di inizio;
					\item durata prevista;
					\item elenco di attività da svolgere.
				\end{itemize}
				A questi incontri partecipano solo i membri del gruppo. Coloro i quali, per diverse motivazioni, non possono essere presenti, sono tenuti a segnalarlo in fase di definizione dell'incontro e sarà compito del responsabile aggiornarli successivamente.
				Gli incontri interni si distinguono in due tipologie:
				\begin{itemize}
					\item \textbf{incontri di persona}: incontri svolti principalmente presso le strutture dell'Università degli Studi di Padova nei quali il gruppo si incontra di persona e discute dell'ordine del giorno prefissato. Questa tipologia viene scelta quando, per impegni universitari, la maggior parte dei membri del gruppo si trova in luoghi vicini oppure per prendere decisioni molto importanti;
					\item \textbf{incontri telematici}: chiamate o video chiamate di gruppo svolte attraverso l'applicazione Skype. Questa tipologia viene scelta quando si necessita di un incontro, ma la maggior parte dei membri del gruppo si trova scomoda o impossibilitata a raggiungere le sedi dell'Università. 
				\end{itemize}
			\subparagraph{Verbali di incontri interni}\mbox{}\\ [1mm]
				Un segretario, nominato dal responsabile, è incaricato della redazione di un \textit{Verbale} della riunione che tiene traccia di:
				\begin{itemize}
					\item \textbf{Luogo};
					\item \textbf{Data};
					\item \textbf{Ora d'inizio};
					\item \textbf{Ora di fine};
					\item \textbf{Partecipanti};
					\item \textbf{Argomenti trattati}: argomenti all'ordine del giorno che sono trattati nell'incontro da verbalizzare;
					\item \textbf{Discussioni e decisioni prese}: discussioni e decisioni prese in merito a tematiche interne al gruppo.
				\end{itemize}
			\subparagraph{Incontri esterni}\mbox{}\\ [1mm]
				Gli incontri esterni vengono concordati se i membri del gruppo, il committente o il proponente lo ritengono necessario. In tal caso, sarà compito del responsabile di progetto\glosp definire data e ora dell'incontro, in accordo con le due parti, attraverso gli appositi canali di comunicazione descritti in precedenza.
				Gli incontri interni si distinguono in due tipologie:
				\begin{itemize}
					\item \textbf{incontri di persona}: incontri svolti di persona presso gli uffici del proponente o del committente, in base a chi è diretto il colloquio, nei quali ognuno espone problemi o dubbi da chiarire e discutendone. Questa tipologia viene scelta per avere dei riscontri immediati qualora tutte le parti fossero disponibili e ci fosse la necessità;
					\item \textbf{incontri telematici}: chiamate o video chiamate attraverso l'applicazione Skype, Hangouts o Zoom sulla base delle richieste del proponente o del committente, in base a chi è diretto il colloquio. Questa tipologia viene scelta quando si necessita di un incontro, ma le parti in causa si trovano impossibilitate ad incontrarsi di persona. 
				\end{itemize}
			\subparagraph{Verbali di incontri esterni}\mbox{}\\ [1mm]
				Un segretario, nominato dal responsabile, è incaricato della redazione di un \textit{Verbale} della riunione tenendo traccia di:
				\begin{itemize}
					\item \textbf{Luogo};
					\item \textbf{Data};
					\item \textbf{Ora d'inizio};
					\item \textbf{Ora di fine};
					\item \textbf{Partecipanti};
					\item \textbf{Argomenti trattati}: argomenti all'ordine del giorno che sono trattati nell'incontro da verbalizzare;
					\item \textbf{Verbale}: esposizione delle discussioni con il proponente e delle decisioni prese in merito.
				\end{itemize}
			\subparagraph{Incontri in situazione eccezionale}\mbox{}\\ [1mm]
			Il Rettore dell'Università degli Studi di Padova ha emesso delle disposizioni in merito alla pandemia in corso causata dal virus Covid-19 nelle quali è ufficializzata la chiusura delle sedi universitarie dal mese di marzo a data da destinarsi. Fino ad allora, tutti gli incontri interni ed esterni sono obbligatoriamente virtuali.
		\paragraph{Gestione degli strumenti di coordinamento}
			\subparagraph{Ticketing}\mbox{}\\ [1mm]
				Il sistema di ticketing permette ai membri del gruppo di gestire le tutte le attività del progetto\glo.
				Per eseguire il ticketing vengono utilizzati gli strumenti dell'Issue Tracking System presenti gratuitamente in GitHub ed il funzionamento è il seguente:
				il responsabile di progetto\glosp ha l'incarico di assegnare i compiti ai membri del gruppo attraverso le issue. Ognuna di esse contiene titolo, descrizione, milestone, le project board Kanban automatizzate associate e il destinatario.
				Le project board Kanban automatizzate permettono al responsabile di progetto\glosp di monitorare l'andamento delle attività tracciando lo stato di ogni singolo compito che può essere:
				\begin{itemize}
					\item To Do (da fare);
					\item In Progress (in lavorazione);
					\item Done (concluso).
				\end{itemize}
				L'associazione con le milestone\glosp definisce le scadenze delle singole issue.
				Inoltre è presente una comoda funzionalità di ricerca per le issue assegnate che ne agevola ulteriormente la gestione.
				
		\subsubsection{Metriche di qualità}
			\paragraph{OP-7 Pianificazione efficace delle risorse}
			\paragraph*{Obiettivo}\mbox{}\\ [1mm]
			L'obiettivo è pianificare le risorse per garantire un corretto avanzamento del progetto\glosp monitorando e rispettando costi e tempistiche.
			\paragraph*{Metriche}\mbox{}\\ [1mm]
			Le metriche\glosp utilizzate sono:
			\begin{itemize}		
				\item \textbf{M07 Planned value}: indica il valore dei lavoro pianificato fino a quel momento;
				\begin{itemize}
					\item[] \textbf{formula}: BAC $\cdot$ \%lavoro pianificato con BAC=budget totale;
				\end{itemize}
				
				\item \textbf{M08 Earned value}: indica il valore del lavoro svolto fino a quel momento;
				\begin{itemize}
					\item[] \textbf{formula}: BAC $\cdot$ \%lavoro svolto BAC=budget totale;
				\end{itemize} 
				
				\item \textbf{M09 Actual cost}: indica il denaro speso fino a quel momento; 
				\begin{itemize}
					\item[] \textbf{formula}: si calcola con un valore intero;
				\end{itemize}
				
				\item \textbf{M10 Cost performance index}: rappresenta un indice per i costi effettivi rispetto a quelli previsti, se è $\le$1 indica che si sta spendendo più di quanto preventivato;
				\begin{itemize}
					\item[] \textbf{formula}: $\frac{EV}{AC}$; 
				\end{itemize}
				
				\item \textbf{M11 Schedule performance index}: rappresenta un indice per i tempi effettivi rispetto a quelli previsti, se è $\le$1 indica che si sta impiegando più tempo di quanto preventivato;
				\begin{itemize}
					\item[] \textbf{formula}: $\frac{EV}{PV}$; 
				\end{itemize}
				
				\item \textbf{M12 Estimated cost at completion}: rappresenta il budget totale stimato del progetto\glosp con il CPI del momento;
				\begin{itemize}
					\item[] \textbf{formula}: $AC+\frac{BT-EV}{CPI}$; 
				\end{itemize} 
				
				\item \textbf{M13 Schedule at completion}: rappresenta il tempo totale stimato del progetto\glosp con SPI del momento;
				\begin{itemize}
					\item[] \textbf{formula}: $\frac{TT}{SPI}$; con TT=Tempo Totale. 
				\end{itemize} 
			\end{itemize}
		
			\paragraph{OP-Q8 Prevenzione dei rischi}
			\paragraph*{Obiettivo}\mbox{}\\ [1mm]
			L'obiettivo è ed analizzare i rischi al fine di trovare metodologie efficaci per preventivarli e successivamente prevenirli.
			\paragraph*{Metriche}\mbox{}\\ [1mm]
			Le metriche\glosp utilizzate sono:
			\begin{itemize}
				\item \textbf{M14 Rischi non preventivati}: indica il numero di rischi non preventivati che si verificano. Un valore elevato indica una superficialità nell'individuazione dei rischi possibili;
				\begin{itemize}
					\item[] \textbf{formula}: si calcola con un valore intero partendo da 0 e andando a fare un incremento ogni volta che si verifica un imprevisto.
				\end{itemize}
			\end{itemize}
			\subsubsection{Strumenti di supporto}
				Gli strumenti utilizzati sono i seguenti:
				\begin{itemize}
					\item \textbf{Telegram}: strumento di messaggistica istantanea utilizzato inizialmente dal gruppo per il primo incontro;
					\item \textbf{Slack}: strumento utilizzato per le comunicazione interne del gruppo;
					\item \textbf{Google Drive}: strumento utilizzato per file che non sono oggetto di cambiamenti frequenti come il \textit{Glossario} o immagini di vario tipo;
					\item \textbf{Google Calendar}: strumento utilizzato per tracciare e semplificare la gestione degli impegni e fornire a tutti notifiche e promemoria per gli stessi;
					\item \textbf{Gmail}: strumento di posta elettronica utilizzato per le comunicazioni esterne.
					\item \textbf{Skype}: strumento utilizzato per effettuare chiamate o video chiamate per alcuni incontri interni o esterni;
					\item \textbf{Hangouts}: strumento utilizzare per effettuare chiamate e video chiamate per alcuni incontri esterni;
					\item \textbf{Zoom}: strumento utilizzato per effettuare chiamate e videochiamate per alcuni incontri esterni; 
					\item \textbf{Git}: sistema di controllo di versionamento\glo;
					\item \textbf{GitKraken, GitHub Desktop}: interfacce l'utilizzo di GitHub in locale;
					\item \textbf{GitHub}: strumento utilizzato per il proprio Issue Tracking System e sistema di versionamento\glosp dei file.
				\end{itemize}
	\subsection{PRC-Q10 Formazione}
		\subsubsection{Scopo}
			Lo scopo di questo processo\glosp è formare ogni membro del gruppo per garantire che ognuno abbia le conoscenze e le competenze adeguate allo svolgimento dei compiti assegnati.
		\subsubsection{Aspettative}
			L'aspettativa del gruppo, nel caso di questo processo\glo, è riuscire a ottimizzare i tempi non impiegando tutti i membri nell'apprendimento di uno stesso strumento o tecnologia, ma lavorare in modo parallelo, al fine di assimilare più concetti possibili per poi condividerli con gli altri componenti.
		\subsubsection{Descrizione}
			Il processo\glosp di formazione è composto da un insieme di attività che definiscono come vengono formati i componenti del gruppo.
			Lo studio necessario per svolgere i compiti e utilizzare gli strumenti viene principalmente eseguito in modo individuale e autonomo da un componente del gruppo definito periodicamente dal responsabile di progetto\glo. Al termine dello studio, egli deve comunicare ciò che ha imparato agli altri membri per garantire che tutti apprendano le nozioni.
			Inoltre, qualora un componente ritenga di non avere ancora sufficienti capacità o conoscenze per svolgere un compito assegnato, si deve rivolgere al responsabile di progetto\glosp che ha l'incarico di organizzare un seminario mirato per il suo apprendimento sia tramite un incontro fisico che in via telematica.
		\subsubsection{Attività}
		\paragraph{Condivisione del materiale}\mbox{}\\ [1mm]
			Il materiale del progetto\glosp è inserito all'interno di un sistema di versionamento\glosp che permette ai membri del gruppo di integrare le proprie conoscenze e apprendere dal lavoro altrui.
			Per ogni documento, il materiale utilizzato per la comprensione dei concetti e la stesura è indicato nei "Riferimenti Informativi".
			Inoltre sono presenti i "Riferimenti Normativi" dentro i quali è presente il documento \textit{Norme di Progetto v. 4.1.1} che tutti devono rispettare.
		\subsubsection{Strumenti di supporto}
		\begin{itemize}
		 	\item \textbf{Skype}: strumento utilizzato per effettuare chiamate o video chiamate per gli incontri di formazione in via telematica;
	 	\end{itemize}
			

	%\pagebreak
	%%\appendix
%\section{Standard di qualità}
%	\subsection{ISO/IEC 25010}
%		Lo standard internazionale ISO/IEC 25010 è utilizzato per misurare la qualità del software.
%		Per farlo utilizza una serie di parametri e caratteristiche che possono essere utilizzate per valutare la qualità del software, grazie alle metriche\glosp fornite dallo standard ISO/IEC 25023.
%		\subsubsection{Metriche di qualità interna}
%			La qualità interna del software è relativa alle attività di progettazione\glosp e codifica, essa rappresenta la qualità della struttura e del codice del prodotto\glo.
%			La misurazione della qualità interna viene effettuata principalmente attraverso l'analisi statica\glosp del codice.
%			Questi controlli sono utili perché semplificano il lavoro dei programmatori e permettono di apportare miglioramenti strutturali al software senza la necessità di eseguirlo.
%			Con un buon livello di qualità interna si ottiene una solida struttura del prodotto\glo, come base per la qualità esterna e in uso.
%		\subsubsection{Metriche di qualità esterna}
%			La qualità esterna del software è relativa alle proprietà dinamiche e comportamentali del software. Essa viene infatti misurata attraverso l'analisi dinamica, ovvero tramite l'esecuzione del codice e lo svolgimento di test.
%		\subsubsection{Metriche di qualità in uso}
%			La qualità in uso è probabilmente la più difficile da misurare poiché è valutata in un ambito di utilizzo reale, con la partecipazione degli utenti.
%			Essa dipende dalla qualità interna e da quella esterna; infatti per avere una buona qualità in uso è necessario un buon livello di qualità interna ed esterna.

	%\pagebreak
	%\subsubsection{Modello di qualità del prodotto}
	\paragraph{Qualità dei documenti}
	Le metriche utilizzate per misurare la qualità dei documenti sono:
	\begin{itemize}
		\item \textbf{Indice di Gulpease}: indica il livello di leggibilità di un testo: più alto questo valore e più è alta la leggibilità del testo. In particolare con valore inferiore a 80 risultata difficile da leggere per chi ha la licenza elementare; inferiore a 60 risultata difficile da leggere per chi ha la licenza media; inferiore a 40 risultata difficile da leggere per chi ha un diploma superiore;
		\begin{itemize}
			\item[] \textbf{formula}: 89+$\frac{300\cdot numero \; di \; frasi-10\cdot numero \; di \; lettere}{numero \; di \; parole}$;
		\end{itemize}                	
		\item \textbf{Correttezza ortografica}: indica la quantità di errori ortografici presenti nel documento.
		\begin{itemize}
			\item[] \textbf{formula}: si calcola con un valore intero;
		\end{itemize}
	\end{itemize}
	\paragraph{Qualità del software} 
	Per descrivere il modello di qualità del prodotto\glosp software il nostro gruppo ha deciso di utilizzare lo standard ISO/IEC 25010, che determina le caratteristiche della qualità di un prodotto\glosp software che verranno valutate. Il modello di qualità del prodotto\glo, definito nello standard, è suddiviso nelle seguenti caratteristiche:
		\subparagraph{Idoneità funzionale} \mbox{}\\[1mm]
		Per adeguatezza funzionale si intende la capacità di un prodotto\glosp software di fornire funzioni che soddisfano i requisiti impliciti ed espliciti, quando usati in un determinato contesto sotto specifiche condizioni. \\
		Le metriche\glosp sono:
		\begin{itemize}
			\item \textbf{M21 Percentuale di requisiti obbligatori soddisfatti}: indica la percentuale di requisiti obbligatori che sono stati soddisfatti. Questo valore deve essere massimizzato per garantire che i vincoli col proponente vengano rispettati;
			\begin{itemize}
				\item[] \textbf{formula}: $\frac{requisiti \; obbligatori \; soddisfatti}{requisiti \; obbligatori \; totali}$;
			\end{itemize} 
			\item \textbf{M22 Percentuale di requisiti desiderabili soddisfatti}:
			indica la percentuale di requisiti desiderabili che sono stati soddisfatti. Un valore elevato denota una maggior copertura delle richieste del proponente;
			\begin{itemize}
				\item[] \textbf{formula}: $\frac{requisiti \; desiderabili \; soddisfatti}{requisiti \; desiderabili \; totali}$.
			\end{itemize} 
			\item \textbf{M23 Percentuale di requisiti opzionali soddisfatti}:
			indica la percentuale di requisiti opzionali che sono stati soddisfatti. Un valore elevato denota una maggior copertura delle richieste non indispensabili  del proponente;
			\begin{itemize}
				\item[] \textbf{formula}: $\frac{requisiti \; opzionali \; soddisfatti}{requisiti \; opzionali \; totali}$.
			\end{itemize} 
			\item \textbf{M24 Percentuale di test passati}:
			indica la percentuale di test passati. Un valore basso indica che la qualità del prodotto è insufficiente. 
			\begin{itemize}
				\item[] \textbf{formula}: $\frac{test \; passati}{totale \; test \; effettuati}$.
			\end{itemize}
		\end{itemize}
	
		\subparagraph{Affidabilità} \mbox{}\\[1mm]
		Un prodotto\glosp software è considerato affidabile quando mantiene un livello di prestazioni costante per un determinato intervallo di tempo sotto precise condizioni che possono variare nel tempo.\\
		Le metriche\glosp sono:
		\begin{itemize}
			\item \textbf{M25 Densità degli errori}: indica la percentuale di errori critici rilevati; 
			\begin{itemize}
				\item[] \textbf{formula}: $\frac{N_{ER}}{N_{T}}\cdot100$ dove $N_{ER}$ sono gli errori rilevati e $N_T$ sono i test eseguiti.
			\end{itemize}
		\end{itemize}
	
		\subparagraph{Efficienza prestazionale} \mbox{}\\
		L'efficienza prestazionale è una caratteristica che indica le prestazioni rispetto al numero di risorse utilizzate sotto condizioni precise. Dato che il nostro progetto si appoggia sul sistema Grafana\glosp e su algoritmi forniti dal proponente, l'efficienza prestazionale viene delegata a questi due strumenti.
		%\begin{itemize}
		%	\item \textbf{Comportamento rispetto al tempo}: la capacità del software di eseguire le funzioni in un tempo di risposta, tempo processionale e un volume di produzione che soddisfa i requisiti;
		%	\item \textbf{Utilizzo delle risorse}: la capacità del software di eseguire le funzioni usando un numero di risorse disponibili che soddisfa i requisiti;
		%	\item \textbf{Capacità}: la capacità del software di avere limiti prestazionali massimi che soddisfano i requisiti.
		%\end{itemize}
		%\paragraph{Compatibilità} \mbox{}\\
		%La compatibilità è la caratteristica di un software di scambiare informazioni con altri prodotti\glo, ed eseguire le sue funzioni mentre viene condiviso lo stesso ambiente hardware o software.
		%\begin{itemize}
		%	\item \textbf{Coesistenza}: la capacità del software di eseguire le sue funzionalità senza impatti negativi mentre condivide lo stesso ambiente con un altro software;
		%	\item \textbf{Interoperabilità}: la capacità di due o più software di scambiare informazioni tra loro e usare tali dati per i loro scopi.
		%\end{itemize}
			
		\subparagraph{Sicurezza} \mbox{}\\
		Per sicurezza in un prodotto\glosp software si intende il grado di protezione che il software ha sui dati che manipola per evitare che persone o servizi non autorizzati li ricevano.
		\begin{itemize}
			\item \textbf{Confidenzialità}: capacità del software di assicurare la confidenzialità dei dati alle persone che hanno l'autorizzazione di accedere a tali informazioni;
			\item \textbf{Integrità}: capacità del software di evitare modifiche o accessi non autorizzati;
			\item \textbf{Non ripudio}: capacità del software di dimostrare che eventi o azioni sono avvenuti, in modo da non ripudiarli in un secondo momento;
			\item \textbf{Responsabilità}: capacità del software di risalire a una sola entità a partire da una azione di una entità;
			\item \textbf{Autenticità}: capacità del software di provare l'identità delle risorse utilizzate.
		\end{itemize}
	
		\subparagraph{Manutenibilità} \mbox{}\\[1mm]
		Per manutenibilità si intende la capacità di un prodotto\glosp software di essere modificato e aggiornato senza l'introduzione di problemi e anomalie in seguito a cambiamenti dell'ambiente e/o dei requisiti.
		Le metriche\glosp sono:
		\begin{itemize}
			\item \textbf{M26 Structural fan-in}: indica quante componenti utilizzano un dato modulo, un valore elevato indica un alto riuso del modulo stesso;
			\begin{itemize}
				\item[] \textbf{formula}: si calcola con un valore intero;
			\end{itemize}
			\item \textbf{M27 Structural fan-out}: indica quante componenti vengono utilizzate dalla componente in esame, un valore elevato indica un alto accoppiamento delle componenti, è quindi preferibile mantenerlo il più basso possibile;
			\begin{itemize}
				\item[] \textbf{formula}: si calcola con un valore intero.
			\end{itemize}
		\end{itemize}
		\item \textbf{M32 Presenza di code smells}: indica la presenza di debolezze nel codice;
		\begin{itemize}
			\item[] \textbf{formula}: si calcola con un valore intero.
		\end{itemize}
	\end{itemize}
	\paragraph{Portabilità} \mbox{}\\ 
	Per portabilità si intende la caratteristica del software di essere utilizzato in modo efficacie quando esso viene trasferito su diversi hardware o software.
	\begin{itemize}
		\item \textbf{Adattabilità}: la capacità del software di essere adattato per hardware o software in continuo cambiamento;
		\item \textbf{Installabilità}: la capacità del software di essere installato o disinstallato efficacemente in un ambiete specifico;
		\item \textbf{Sostituibilità}: la capacità del software di rimpiazzare un altro prodotto\glosp con lo stesso scopo nello stesso ambiente.
	\end{itemize}

	\paragraph{Usabilità} \mbox{}\\[1mm]
	L'usabilità è una caratteristica di un software definita come l'efficacia, l'efficienza e la soddisfazione con le quali determinati utenti raggiungono determinati obiettivi in determinati contesti.

	% ...
	%\input{res/inserire nome sezione n} 
	
	
	
	%Input del file "decision_table" della cartella locale "res/"
	%% \textbf = grassetto; \Large = font più grande
% \rowcolors{quanti colori alternare}{colore numero riga pari}{colore numero riga dispari}: colori alternati per riga
% \rowcolor{color}: cambia colore di una riga
% p{larghezza colonna}: p è un tipo di colonna di testo verticalmente allineata sopra, ci sarebbe anche m che è centrata a metà ma non è precisa per questo utilizzo TBStrut; la sintassi >{\centering} indica che il contenuto della colonna dovrà essere centrato
% \TBstrut fa parte di alcuni comandi che ho inserito in config.tex che permetto di aggiungere un po' di padding al testo
% \\ [2mm] : questra scrittura indica che lo spazio dopo una break line deve essere di 2mm

%\setcounter{secnumdepth}{0}
%\hfill \break
%\textbf{\Large{Diario delle modifiche}} \\
\section{Riepilogo tracciamenti}
\rowcolors{2}{gray!25}{gray!15}
\begin{longtable} {
		>{\centering}p{17mm} 
		%>{\centering}p{19.5mm}
		%>{\centering}p{24mm} 
		%>{\centering}p{24mm} 
		>{}p{120mm}}
	\rowcolor{gray!50}
	\textbf{Codice} & \multicolumn{1}{c}{\textbf{Decisione}} \\%\textbf{Decisione} \\ %\TBstrut \\
	VI\_2.1 & Divisione dei ruoli per inizio stesura documenti. \TBstrut \\ [2mm]
	VI\_2.2 & Studio preliminare e divisione del carico di lavoro per le \textit{Norme di Progetto v. 1.1.1} \TBstrut \\ [2mm]
	VI\_2.3 & Stesura domande per l'incontro con il proponente \textit{Zucchetti}. \TBstrut \\ [2mm]
	
\end{longtable} %Tabella delle decisioni (solo per i verbali)
\end{document}